\chapter{The Modular Trace}
\label{ch:modular-trace}

A chiral algebra carries a modular characteristic $\kappa_{\mathrm{ch}}$; a Calabi--Yau category carries a categorical trace $\Tr \colon \HH_d(\cC) \to k$; a Calabi--Yau manifold carries a holomorphic Euler characteristic $\chi(\cO_X)$. The tempting shortcut identifying $\kappa_{\mathrm{ch}}$ with the topological Euler characteristic $\chi_{\mathrm{top}}/24$ is \emph{wrong in every computed case}. The right target is $\chi(\cO_X)$, and even this equals $\kappa_{\mathrm{ch}}$ only on the restricted slice $(d=2,\ h^{1,0}=0)$.

For the elliptic curve, $\chi_{\mathrm{top}} = 0$ but $\kappa_{\mathrm{ch}}(H_1) = 1$. For $K3$, $\chi_{\mathrm{top}}/24 = 1$ but $\kappa_{\mathrm{ch}}(\cA_{K3}) = 2 = \dim_\C$. For $K3 \times E$, two different modular characteristics appear: $\kappa_{\mathrm{ch}}^{\mathrm{Heis}} = 3$ from the Heisenberg-level chiral de Rham complex (cf.\ Remark~\ref{rem:beauville-kappa-formula-subscript-split}) and $\kappa_{\mathrm{BKM}} = 5$ from the Borcherds lift weight. For the resolved conifold, $\chi_{\mathrm{top}}/24 = 1/12$ but $\kappa_{\mathrm{ch}} = 1$. The topological invariant is not what the chiral algebra sees.

This chapter replaces the wrong identification by the right one. The CY trace, properly refined to negative cyclic homology $\HC^-_d(\cC)$, determines $\kappa_{\mathrm{ch}}(A_\cC)$ through the CY-to-chiral correspondence programme $\{\Phi_d\}$. At $d = 2$, the resulting equality $\kappa_{\mathrm{ch}}(A_\cC) = \chi^{\CY}(\cC) = \chi(\cO_X)$ is proved (Proposition~\ref{prop:cy-kappa-d2}). At $d \neq 2$, $\kappa_{\mathrm{ch}}$ and $\chi(\cO_X)$ differ in general: the Heisenberg-level characteristic $\kappa_{\mathrm{ch}}^{\mathrm{Heis}}$ is additive under products while $\chi(\cO_X)$ is multiplicative (K\"unneth), so $\kappa_{\mathrm{ch}}^{\mathrm{Heis}}(K3 \times E) = 3 \neq 0 = \chi(\cO_{K3 \times E})$ (cf.\ Remark~\ref{rem:beauville-kappa-formula-subscript-split}). The corrected $d = 3$ formula is dimension-stratified (Conjecture~\ref{conj:cy-kappa-identification}). The genus-$g$ obstruction tower $\mathrm{obs}_g(A_\cC) = \kappa_{\mathrm{ch}}(A_\cC) \cdot \lambda_g$ then encodes the higher-genus CY invariants on the uniform-weight lane, with the multi-weight cross-channel correction $\delta F_g^{\mathrm{cross}}$ from Vol~I appearing at $g \geq 2$ for families with fields of distinct conformal weights.

\section{CY trace as modular characteristic}
\label{sec:cy-trace-kappa}

%: kappa always subscripted in Vol III.
The CY trace $\Tr \colon \HH_d(\cC) \to k$ determines the modular characteristic $\kappa_{\mathrm{ch}}(A_\cC)$.

\begin{theorem}[CY modular characteristic: Theorem CY-D]
\label{thm:cy-modular-characteristic}
\ClaimStatusProvedHere{}
Fix the Vol~III convention that the operational chiral modular
characteristic is the K\"unneth-additive invariant
\[
  \kappa_{\mathrm{ch}}(X)=\kappa_{\mathrm{ch}}^{\mathrm{K}}(X),
\]
while the Hodge comparator is written separately as
\[
  \kappa_{\mathrm{ch}}^{\mathrm{Hodge}}(X)
  := \sum_{q=0}^{d} (-1)^q h^{0,q}(X)
  = \chi(\cO_X).
\]
Let $\cC$ be a smooth proper CY category and let
$A_\cC=\Phi(\cC)$.
\begin{enumerate}[label=(\roman*)]
 \item If $d=2$ and $\HH_{-1}(\cC)=0$ \textup{(}equivalently, for
       $\cC=D^b\Coh(X)$, $h^{1,0}(X)=0$\textup{)}, then
       \[
         \kappa_{\mathrm{ch}}(A_\cC)
         = \chi^{\CY}(\cC)
         = \chi(\cO_X).
       \]
       This is the proved CY-D$_2$ slice: the Serre-duality parity
       argument kills the one-loop correction.
 \item For compact CY input of dimension $d \ge 3$, the comparison
       between $\kappa_{\mathrm{ch}}$ and
       $\kappa_{\mathrm{ch}}^{\mathrm{Hodge}}$ is stratified by the
       Beauville--Bogomolov holonomy type:
       \begin{enumerate}[label=\textup{(\alph*)}]
         \item if $d$ is odd, then
               $\kappa_{\mathrm{ch}}^{\mathrm{Hodge}}(X)=0$ by Serre
               cancellation, but bare $\kappa_{\mathrm{ch}}$ need not
               vanish;
         \item if $d$ is even and $X$ is strict CY, then the Hodge
               comparator is the genuine Serre-duality endpoint
               contribution
               $\kappa_{\mathrm{ch}}^{\mathrm{Hodge}}(X)=2$;
         \item if $d=2n$ and $X$ is irreducible
               holomorphic-symplectic, then
               $\kappa_{\mathrm{ch}}^{\mathrm{Hodge}}(X)=n+1$,
               generated by $1,\sigma,\dots,\sigma^n$.
       \end{enumerate}
 \item The product $K3 \times E$ is the basic obstruction to the naive
       identity $\kappa_{\mathrm{ch}}=\chi(\cO_X)$:
       \[
         \kappa_{\mathrm{ch}}^{\mathrm{Hodge}}(K3 \times E)
         = \chi(\cO_{K3 \times E}) = 0,
         \qquad
         \kappa_{\mathrm{ch}}(K3 \times E)=3.
       \]
       The left-hand side is K\"unneth-multiplicative; the right-hand
       side is K\"unneth-additive.
 \item On the uniform-weight lane, the genus-$g$ obstruction satisfies
       $\mathrm{obs}_g(A_\cC)=\kappa_{\mathrm{ch}}(A_\cC)\cdot\lambda_g$
       unconditionally at genus~$1$ and at higher genus up to the usual
       multi-weight cross-channel correction $\delta F_g^{\mathrm{cross}}$.
       At genus $g\ge 2$ for multi-weight algebras, the scalar formula
       receives the nonvanishing cross-channel correction (Vol~I
       op:multi-generator-universality, resolved negatively;
       $\delta F_2(\mathcal{W}_3)=(c{+}204)/(16c)>0$). The shadow
       obstruction tower of $A_\cC$ encodes the higher-genus CY
       invariants of $\cC$.
\end{enumerate}
\end{theorem}

\section{CY Euler characteristic versus modular characteristic}
\label{sec:chi-vs-kappa}

The modular characteristic $\kappa_{\mathrm{ch}}$ is an invariant of the quantum chiral algebra $A_\cC$, not of the underlying topology. It differs from the topological Euler characteristic $\chi$ in every known case.

\begin{proposition}[$\chi/24 \neq \kappa_{\mathrm{ch}}$ in general]
\label{prop:chi-kappa-discrepancy}
\ClaimStatusProvedHere
The ratio $\chi_{\mathrm{top}}/24$ and the modular characteristic $\kappa_{\mathrm{ch}}$ disagree for most CY threefolds. The following table records all cases where both are known.
%: The column ``$\kappa_{\mathrm{ch}}$'' records the modular
% characteristic of the chiral algebra output. For K3 x E, two
% values appear: kappa_ch = 3 (chiral de Rham, additive under
% products) and kappa_BKM = 5 (Borcherds weight, conjectural as a
% Vol~I modular characteristic). See Remark rem:cy3-kappa-polysemy.
\begin{center}
\begin{tabular}{lcccl}
\toprule
$X$ & $\chi_{\mathrm{top}}(X)$ & $\chi_{\mathrm{top}}/24$ & $\kappa_{\mathrm{ch}}$ & Source \\
\midrule
Elliptic curve $E$ & $0$ & $0$ & $1$ & $\kappa_{\mathrm{ch}}(H_1) = 1$ \\
$K3$ surface & $24$ & $1$ & $2$ & $\kappa_{\mathrm{ch}}(\cA_{K3}) = 2$ (Prop.~\ref{prop:kappa-k3}) \\
$K3 \times E$ & $0$ & $0$ & $3$ / $5^*$ & $\kappa_{\mathrm{ch}} = 3$ (proved); $\kappa_{\mathrm{BKM}} = 5$ (conj.) \\
Conifold & $2$ & $1/12$ & $1$ & Resolved conifold \\
\bottomrule
\end{tabular}
\end{center}
\noindent{}${}^*$For $K3 \times E$, the chiral de Rham complex gives $\kappa_{\mathrm{ch}} = 3 = \dim_\C$; the Borcherds lift weight gives $\kappa_{\mathrm{BKM}} = 5 = \mathrm{wt}(\Delta_5)$. These are modular characteristics of \emph{different} algebras. The chiral algebra $A_{K3 \times E}$ is the output of Theorem~CY-A$_3$ (Theorem~\ref{thm:cy-to-chiral-d3}) applied to $D^b\Coh(K3 \times E)$; the identification $\kappa_{\mathrm{BKM}} = 5$ as a Vol~I modular characteristic is conditional on the CY-D programme (motivic DT identification).

\smallskip\noindent In particular: $\chi_{\mathrm{top}}/24 \neq \kappa_{\mathrm{ch}}$ in every case.
\end{proposition}

\begin{proof}
Each entry is computed independently. For $E$: the quantum chiral algebra is the Heisenberg $H_1$ with $\kappa_{\mathrm{ch}} = 1$ (the level), while $\chi_{\mathrm{top}}(E) = 0$. For $K3$: the output of $\Phi$ at the categorical input $D^b(\Coh(K3))$ is the Mukai Heisenberg $H_{\mathrm{Muk}}$ (Theorem~\ref{thm:phi-k3-explicit}), with $\kappa_{\mathrm{ch}}(H_{\mathrm{Muk}}) = \chi(\cO_{K3}) = 2$ by Theorem~\ref{thm:cy-modular-characteristic}, while $\chi_{\mathrm{top}}(K3)/24 = 1$. The $\cN=4$ superconformal algebra of the K3 sigma model is a distinct chiral algebra that happens to share $\kappa_{\mathrm{ch}} = 2$ but does not arise as $\Phi$ applied to any categorical input in Vol~III. For $K3 \times E$: $\chi_{\mathrm{top}}(K3 \times E) = \chi(K3) \cdot \chi(E) = 24 \cdot 0 = 0$; the chiral de Rham complex has $\kappa_{\mathrm{ch}} = \kappa_{\mathrm{ch}}(K3) + \kappa_{\mathrm{ch}}(E) = 2 + 1 = 3$ (proved by additivity); the BKM automorphic weight is the distinct quantity $\kappa_{\mathrm{BKM}} = \mathrm{wt}(\Delta_5) = 5$ (see the $\kappa_{\mathrm{ch}}$-spectrum, Example~\ref{ex:kappa-spectrum-k3xe}). For the conifold: the resolved conifold has $\chi_{\mathrm{top}} = 2$ (the total space of $\cO(-1) \oplus \cO(-1) \to \bP^1$ deformation retracts onto the zero section $\bP^1$, so $\chi_{\mathrm{top}} = \chi(\bP^1) = 2$), giving $\chi_{\mathrm{top}}/24 = 1/12$, while $\kappa_{\mathrm{ch}} = 1$.

See \texttt{compute/lib/cy\_euler.py} ($86$ tests) and \texttt{compute/lib/modular\_cy\_characteristic.py} ($80$ tests).
\end{proof}

\begin{remark}[Categorical versus topological Euler characteristic]
\label{rem:categorical-vs-topological-chi}
The CY Euler characteristic $\chi^{\CY}(\cC)$ appearing in Theorem~\ref{thm:cy-modular-characteristic} is the categorical trace $\Tr_{\HH_d(\cC)}(\mathrm{id})$, which is in general \emph{not} the topological Euler characteristic $\chi_{\mathrm{top}}(X)$. The former depends on the CY category $\cC$ (and in particular on the complex structure and derived category), while the latter depends only on the underlying topological manifold.
\end{remark}

\begin{proposition}[Non-multiplicativity of $\kappa_{\mathrm{BKM}}$]
\label{prop:kappa-non-multiplicative}
\ClaimStatusProvedHere
The BKM modular weight $\kappa_{\mathrm{BKM}}$ obeys neither the additive rule that governs $\kappa_{\mathrm{ch}}^{\mathrm{Heis}}$ (Heisenberg-level route) nor the multiplicative rule that governs $\kappa_{\mathrm{cat}}$ (categorical Euler-characteristic route) under products:
\[
 \kappa_{\mathrm{ch}}^{\mathrm{Heis}}(K3) + \kappa_{\mathrm{ch}}^{\mathrm{Heis}}(E) = 2 + 1 = 3, \qquad
 \kappa_{\mathrm{cat}}(K3) \cdot \kappa_{\mathrm{cat}}(E) = 2 \cdot 0 = 0, \qquad
 \kappa_{\mathrm{BKM}}(K3 \times E) = 5.
\]
% WARNING: an earlier draft of this display wrote
%   "kappa_ch(K3) * kappa_ch(E) = 2*1 = 2"
% which mixed Route B (Heisenberg-additive) values under a Route A
% (Hodge-supertrace multiplicative) operation. Route discipline: see
% cy_d_kappa_stratification.tex:411-426.
The discrepancy $5 \neq 3$ and $5 \neq 0$ reflects the fact that the BKM superalgebra $\mathfrak{g}_{\Delta_5}$ for $K3 \times E$ is \emph{not} the tensor product of the $K3$ and $E$ chiral algebras, nor is its weight extracted from a categorical trace. The weight $5 = \mathrm{wt}(\Delta_5)$ is determined by the Borcherds product structure. By contrast, the chiral de Rham modular characteristic $\kappa_{\mathrm{ch}}^{\mathrm{Heis}}(K3 \times E) = \kappa_{\mathrm{ch}}^{\mathrm{Heis}}(K3) + \kappa_{\mathrm{ch}}^{\mathrm{Heis}}(E) = 2 + 1 = 3$ \emph{is} additive in the Heisenberg-level sense (Proposition~\ref{prop:kappa-k3}), while the categorical Euler characteristic $\kappa_{\mathrm{cat}}(K3 \times E) = \chi(\cO_{K3}) \cdot \chi(\cO_E) = 2 \cdot 0 = 0$ is K\"unneth-multiplicative.
\end{proposition}


\begin{proposition}[$\chi(\cO_X) = 0$ for CY manifolds of odd dimension]
\label{prop:chi-O-vanishes-odd-d}
\ClaimStatusProvedHere
For a compact CY$_d$ manifold $X$ with $d$ odd and $\omega_X \simeq \cO_X$, the holomorphic Euler characteristic vanishes:
\[
 \chi(\cO_X) \;=\; \sum_{q=0}^{d} (-1)^q\, h^{0,q}(X) \;=\; 0.
\]
In particular, $\chi(\cO_E) = 0$ for every elliptic curve $E$, and $\chi(\cO_X) = 0$ for every compact CY$_3$ threefold $X$.
\end{proposition}

\begin{proof}
By Serre duality with $\omega_X \simeq \cO_X$: $H^q(X, \cO_X) \cong H^{d-q}(X, \cO_X)^\vee$, so $h^{0,q} = h^{0,d-q}$. When $d$ is odd, every index $q$ is paired with $d-q \neq q$ (since $d$ is odd, $q$ and $d-q$ have opposite parities), and $(-1)^q h^{0,q} + (-1)^{d-q} h^{0,d-q} = (-1)^q h^{0,q} (1 + (-1)^d) = 0$ (since $(-1)^d = -1$ when $d$ is odd). The sum telescopes to zero.

\medskip\noindent\textit{Verification}: 76 tests in \texttt{compute/lib/cy\_d\_kappa\_d3.py} covering the Serre vanishing for 5 standard CY manifolds.
\end{proof}

\begin{remark}[Consequence for CY-D at $d = 3$]
\label{rem:chi-O-obstruction-cy-d}
Proposition~\textup{\ref{prop:chi-O-vanishes-odd-d}} implies that the identification $\kappa_{\mathrm{ch}} = \chi(\cO_X)$ is \emph{automatically false} at $d = 3$ whenever $\kappa_{\mathrm{ch}} \neq 0$. Since $\kappa_{\mathrm{ch}}^{\mathrm{Heis}}(K3 \times E) = 3 \neq 0$ (Proposition~\textup{\ref{prop:categorical-euler}}; Heisenberg-level branch, cf.\ Remark~\ref{rem:beauville-kappa-formula-subscript-split}) and $\chi(\cO_{K3 \times E}) = 0$, the formula $\kappa_{\mathrm{ch}} = \chi(\cO_X)$ fails. The correct CY-D conjecture at $d \geq 3$ (Conjecture~\textup{\ref{conj:cy-kappa-identification}}) uses the categorical CY Euler characteristic $\chi^{\CY}(\cC)$, a quantity that in general differs from $\chi(\cO_X)$ and includes the quantum correction from the $\bS^d$-framing.
\end{remark}


\section{Modularity constraints on the shadow tower}
\label{sec:modularity-constraints}

When a BKM denominator identity controls the shadow obstruction tower of a CY$_3$ category $\cC$, the automorphic form imposes rigid constraints on the tower amplitudes.

\begin{theorem}[BKM modularity constraints]
\label{thm:bkm-modularity-constraints}
\ClaimStatusProvedHere
Let $\cC$ be a CY$_3$ category whose DT partition function $Z^{\DT}(\cC)$ is controlled by a Siegel modular form $\Phi$ of weight $\kappa_{\mathrm{BKM}}$ for an arithmetic group $\Gamma \subset \Sp_4(\Q)$. Then the shadow tower amplitudes $\{F_g\}$ satisfy the following constraints:
\begin{enumerate}[label=(\roman*)]
 \item \emph{Integrality}: $\kappa_{\mathrm{BKM}} \in \Z_{>0}$, since $\kappa_{\mathrm{BKM}}$ is the weight of a Siegel modular form. In particular, the fractional values $\kappa_{\mathrm{ch}} = -25/3$ (quintic threefold, $\chi = -200$) and $\kappa_{\mathrm{ch}} = -1/6$ (banana configuration) are ruled out from admitting BKM lifts.
 \item \emph{Fourier--Jacobi structure}: $\Phi(\tau, z, \omega) = \sum_{m \geq 0} \phi_m(\tau, z) \, e^{2\pi i m \omega}$, where each $\phi_m$ is a Jacobi form of weight $\kappa_{\mathrm{BKM}}$ and index $m$. The shadow tower degree-$r$ contribution maps to the Fourier--Jacobi expansion:
 \begin{align*}
 \text{degree } 2 \;(\kappa_{\mathrm{ch}}) &\;\longleftrightarrow\; \phi_0 \quad \text{(constant term, Eisenstein part)}, \\
 \text{degree } 3 \;(\text{cubic shadow}) &\;\longleftrightarrow\; \phi_1 \quad \text{(first FJ coefficient)}.
 \end{align*}
 \item \emph{Discriminant-degree correspondence}: for the Jacobi form $\phi_{0,1}$ controlling the Borcherds lift, the discriminant $D = r^2 - 4nm$ of a Fourier coefficient $c(n, r)$ classifies the shadow degree:
 \begin{alignat*}{2}
 D < 0 &\quad\text{(real)}: &\quad& \text{degree } 2 \text{ (Kac--Moody roots)}, \\
 D = 0 &\quad\text{(lightlike)}: &\quad& \text{degree } 3 \text{ (affine roots)}, \\
 D > 0 &\quad\text{(timelike)}: &\quad& \text{degree } \geq 4 \text{ (imaginary roots)}.
 \end{alignat*}
\end{enumerate}
\end{theorem}

\begin{proof}
Part~(i): the weight of a Siegel modular form for $\Sp_4(\Z)$ (or a congruence subgroup) is a non-negative integer. The values $\kappa_{\mathrm{ch}} = -25/3$ and $\kappa_{\mathrm{ch}} = -1/6$ are not integers, hence no Siegel modular form of that weight exists.

Part~(ii): the Fourier--Jacobi expansion is standard (Eichler--Zagier). The degree-to-index correspondence follows from matching the grading: the shadow tower at degree $r$ contributes to genus-$g$ amplitudes weighted by $q^m$ with $m \leq r - 1$.

Part~(iii): the discriminant classification follows from the root-type decomposition of the BKM superalgebra: real roots ($D < 0$) correspond to simple-pole OPE singularities (Kac--Moody type, degree $2$), lightlike roots ($D = 0$) to affine null-root contributions (degree $3$), and imaginary roots ($D > 0$) to higher-multiplicity contributions (degree $\geq 4$). The multiplicity $\mult(D)$ of imaginary roots at discriminant $D > 0$ encodes the quartic and higher shadow corrections.

See \texttt{compute/lib/cy3\_modularity\_constraints.py} ($111$ tests).
\end{proof}

\begin{corollary}[Structural constraints for $K3 \times E$]
\label{cor:k3e-structural-constraints}
\ClaimStatusProvedHere
For $K3 \times E$ with $\kappa_{\mathrm{BKM}} = 5$ and controlling form $\Delta_5$ — the weight-$5$ Gritsenko--Nikulin Borcherds product for the Mukai lattice, realised as a cusp form on a paramodular subgroup $\Gamma_{\mathrm{para}} \subset \Sp_4(\Q)$ via the accidental isomorphism $\mathrm{O}^+(2,3) \simeq \PGSp_4$ (see Chapter~\ref{ch:k3-times-e}) — the BKM shadow tower satisfies seven structural constraints:
\begin{enumerate}[label=(\arabic*)]
 \item $\kappa_{\mathrm{BKM}} = 5 \in \Z_{>0}$;
 \item the tower amplitudes $F_g$ lie in the Fourier--Jacobi ring of $\Gamma_{\mathrm{para}}$;
 \item the Hecke eigenvalues of $\Delta_5$ constrain $F_g$ multiplicatively;
 \item the functional equation of $\Delta_5$ imposes a symmetry on the tower;
 \item the Borcherds product representation constrains the root multiplicities;
 \item the Petersson norm of $\Delta_5$ bounds the growth rate of $F_g$;
 \item the theta correspondence $\phi_{0,1} \mapsto \Delta_5$ factors the tower through the Jacobi form $\phi_{0,1}$.
\end{enumerate}
\end{corollary}


\section{Genus expansion and Gromov--Witten invariants}
\label{sec:genus-gw}

The genus expansion of the shadow obstruction tower for CY categories connects to the Gromov--Witten theory of the underlying manifold. For a CY$_3$ manifold $X$ with chiral algebra $A_X = \Phi(\cC_X)$, the genus-$g$ amplitude $F_g(A_X)$ on the uniform-weight lane equals $\kappa_{\mathrm{ch}}(A_X) \cdot \lambda_g$ (Theorem~\ref{thm:cy-modular-characteristic}(ii)), where $\lambda_g = c_g(\mathbb{E})$ is the top Chern class of the Hodge bundle on $\overline{\mathcal{M}}_g$ (the Faber--Pandharipande universal Hodge integral); the identification of this amplitude with the genus-$g$ Gromov--Witten free energy $F_g^{\mathrm{GW}}(X)$ requires the comparison between the shadow tower and the B-model topological string. By Theorem~CY-A$_3$ (Theorem~\ref{thm:cy-to-chiral-d3}), the chiral algebra $A_X = \Phi_3(D^b\Coh(X))$ exists. At $g = 1$, the comparison reduces to the BCOV formula $F_1 = \kappa_{\mathrm{ch}}/24$, unconditionally proved for all families via Vol~I Theorem~D. At $g \geq 2$, the comparison is programme-level: it requires the identification of the shadow tower with the topological string genus expansion. CY-A$_3$ settles the existence of $A_X$; the genus-expansion comparison is a separate step, open at $g \geq 2$.

\section{The CY shadow obstruction tower}
\label{sec:cy-shadow-tower}

The shadow obstruction tower of a CY chiral algebra $A_\cC = \Phi(\cC)$ is the Vol~I datum $\Theta_{A_\cC} \in \mathfrak{g}^{\mathrm{mod}}_{A_\cC}$, with projections $\kappa_{\mathrm{ch}}, C, Q, \ldots$ recording the degree-$2$, degree-$3$, degree-$4$, and higher shadow invariants. For CY$_2$ input (K3, abelian surface), the tower is computable: $\kappa_{\mathrm{ch}} = \chi^{\CY}(\cC) = \chi(\cO_X)$ at $d = 2$ (Theorem~\ref{thm:cy-modular-characteristic}), and the shadow class is determined by the Vol~I G/L/C/M classification applied to $A_\cC$. For CY$_3$ input, the tower is accessible via CY-A$_3$ (Theorem~\ref{thm:cy-to-chiral-d3}); the BKM modularity constraints of Section~\ref{sec:modularity-constraints} provide structural predictions for the tower amplitudes when a Siegel modular form controls the DT partition function.

\section{Complementarity for CY pairs}
\label{sec:cy-complementarity}

Koszul complementarity for CY chiral algebras takes the form $\kappa_{\mathrm{ch}}(A_\cC) + \kappa_{\mathrm{ch}}(A_\cC^!) = K_\cC$, where $K_\cC$ is the family-dependent Koszul conductor (Vol~I). For CY input producing Kac--Moody output, $K_\cC = 0$; for Virasoro-type output, $K_\cC = 13$ (the Virasoro self-dual point $c = 13$). The CY Langlands involution $\cC \mapsto \cC^L$ of Chapter~\ref{ch:geometric-langlands} conjecturally intertwines with Koszul complementarity: $\Phi(\cC^L) \simeq \Phi(\cC)^!$ would identify the Langlands dual category with the Koszul dual chiral algebra (Conjecture~\ref{conj:cy-langlands}(i)).

\smallskip
\noindent\textit{Into Part~\ref{part:connections}.} The Koszul conductor $K_\cC$ is the scalar projection of a richer object: the $r$-matrix $r_{\CY}(z) \in A_\cC \otimes A_\cC((z))$ that encodes the full spectral-parameter braiding. Part~\ref{part:connections} presents seven independent constructions of $r_{\CY}(z)$ (Yangian coproduct twist, Drinfeld-centre half-braiding, KZ monodromy, BKM vertex-operator exchange, MO stable envelope, $\Ainf$-coassociator, Costello $5d$ hCS tree amplitude), all agreeing on $\Phi(\cC)$; the modular trace constructed in this chapter is the scalar shadow of that heptagon.

\section{Modular trace on $\mathbf H_{\Delta_5}$}
\label{sec:modtr-wave13-DNA}

\begin{remark}[Modular trace on $\mathbf H_{\Delta_5}$]
\label{rem:modtr-DNA-1}
The modular trace on $\mathbf H_{\Delta_5}$, defined as
\[
\mathrm{Tr}^{\mathrm{mod}}(\mathbf H_{\Delta_5})(Z) \;=\; \mathrm{Tr}\bigl(q^{L_0 - c/24}\,\zeta^{J_0}\,p^{L_0 - c/24}\bigr),
\]
with $Z = (\tau, z, \tau')\in\mathfrak H_2$ parameterising the Siegel
upper half-space, evaluates to the Gritsenko-Nikulin Igusa cusp form
$\Phi_{10}(Z)$ (via the denominator identity). The three-factor
structure of the trace reflects:
\begin{itemize}
\item $q^{L_0}$: the $\tau$-character of the vertex-operator-sector;
\item $\zeta^{J_0}$: the flavour character of the $\mathfrak{su}(2)^{24}$
      flavour symmetry;
\item $p^{L_0'}$: the $\tau'$-character of the Siegel-modular dual.
\end{itemize}
Primary-lit: Kac 1998 modular invariance, Gritsenko-Nikulin 1996-98,
Eguchi-Ooguri-Tachikawa 2011.
\end{remark}

\section{The modular $S$-matrix as Fricke involution}
\label{sec:modtr-wave16-DNA}

The modular trace of Remark~\ref{rem:modtr-DNA-1} on the Siegel
upper-half space $\mathfrak H_2$ carries a modular $S$-transformation
action whose structure is pinned by the arithmetic of the $\mu_8$-gerbe
on $\overline{\mathcal A_2}$ (Chapter~\ref{ch:k3-chiral-bialgebra-platonic}
Remark~\ref{rem:cydkappa-wave15-gerbe}). At the
$\ell = 8$ Lusztig level of the BKM small-form
(Chapter~\ref{ch:quantum-groups}
Theorem~\ref{thm:qgf-wave16-BKM-MTC}), this $S$-transformation acquires
a geometric identification as the Siegel Fricke involution.

\begin{theorem}[Modular $S$-matrix as Fricke involution on $\mathfrak H_2$]
\label{thm:modtr-wave16-S-matrix-Fricke}
\ClaimStatusProvedHere
Let $\mathcal{MTC}_{\zeta_8}(\Delta_5)$ denote the semisimple
MTC of Chapter~\ref{ch:quantum-groups}
Theorem~\ref{thm:qgf-wave16-BKM-MTC}.  Its modular $S$-matrix
$S_{ij} = \chi_i(-Z^{-1})/\chi_i(Z)$ (CFT normalisation) on the
$M_{24}$-equivariant fusion basis of
Chapter~\ref{ch:k3-chiral-bialgebra-platonic}
Remark~\ref{rem:cydkappa-wave16-fusion-ring} acts on the Siegel
variable $Z\in\mathfrak H_2$ as the Fricke involution
\[
  w_8\colon Z \longmapsto -\bigl(8\,Z\bigr)^{-1},
\]
with fixed locus
\[
  \mathrm{Fix}(w_8)
  \;=\;
  \bigl\{Z\in\mathfrak H_2 \mid 8\,Z\cdot Z = -\mathrm{id}_{2\times 2}\bigr\}
  \;=\;
  H_1\cup H_4,
\]
the union of the two Humbert surfaces classifying genus-$2$ Jacobians
with real multiplication by $\mathbb Z[\sqrt{1}]$ and
$\mathbb Z[\sqrt{4}]=\mathbb Z$-reducible pairs respectively.  The
fixed locus coincides with the $\mu_8$-gerbe-obstructed divisor of
$\overline{\mathcal A_2}$
(Chapter~\ref{ch:k3-chiral-bialgebra-platonic}
Remark~\ref{rem:cydkappa-wave15-gerbe}).
\end{theorem}

\begin{proof}
The CFT-normalised $S$-matrix of a ribbon category $\mathcal C$ is
defined as the Lyubashenko trace
$S_{ij} = \mathrm{tr}_{L_i\otimes L_j}(c_{L_j,L_i}\circ c_{L_i,L_j})$
where $c$ is the braiding (Lyubashenko--Majid 1994 \emph{J.~Algebra}~166
Prop.~3.4).  For $\mathcal{MTC}_{\zeta_8}(\Delta_5)$, the braiding
inherits the Siegel-elliptic dynamical $R$-matrix $R_{\mathrm{Sieg,dyn}}$
of $\mathbf H_{\Delta_5}$; the Siegel-dynamical variable is the modulus
$Z\in\mathfrak H_2$.  Under modular transformation
$Z\mapsto (AZ+B)(CZ+D)^{-1}$ with $\gamma=\begin{psmallmatrix}A&B\\C&D\end{psmallmatrix}
\in\mathrm{Sp}_4(\mathbb Z)$, the $R$-matrix transforms by an
automorphy factor of weight $-\mathrm{wt}(\Delta_5) = -5$ and
multiplier $v_{\Delta_5}$
(Gritsenko--Nikulin 1998 \emph{Amer.~J.~Math.}~120 Lemma~2.3).  The
$S$-matrix, extracted as the trace of the double braiding on
fundamental representations, transforms by the reciprocal automorphy
factor at $S = \begin{psmallmatrix}0&-\mathrm{id}\\\mathrm{id}&0\end{psmallmatrix}$,
reducing to the Fricke involution $w_N\colon Z\mapsto -(NZ)^{-1}$ at
the distinguished level $N = \ell = 8$
(Atkin--Lehner 1970 \emph{Math.~Ann.}~185 eq.~1.7, extended to the
paramodular cover by Gritsenko 1995 \emph{St.~Petersburg Math.~J.}~6
\S3).

\emph{Fixed-locus computation.} $Z\in\mathrm{Fix}(w_8)$ iff
$-(8Z)^{-1} = Z$, i.e.\ $8 Z^2 = -\mathrm{id}$, equivalently
$Z^2 = -\tfrac{1}{8}\mathrm{id}$.  Writing
$Z = \begin{psmallmatrix}\tau_1 & z\\ z & \tau_2\end{psmallmatrix}$,
the condition $\det Z \cdot Z^{-1} = -8 Z$ gives
$\tau_1\tau_2 - z^2 = \pm i/\sqrt{8}$, parameterising a complex
$2$-cycle.  The Humbert surface $H_D\subset\mathfrak H_2$ of
discriminant $D$ (classifying genus-$2$ Abelian surfaces with real
multiplication by an order of discriminant $D$) is defined by
$a\tau_1 + b\tau_2 + c z + d(\tau_1\tau_2 - z^2) + e = 0$ for integers
$(a,b,c,d,e)$ with $D = c^2 - 4(ab+de)$
(Humbert 1899 \emph{Liouville~J.~Math.}~5; van der Geer 1988
\emph{Hilbert Modular Surfaces} Ch.~IX).  At $D = 1$ (reducible
Jacobians $E_1\times E_2$ with common isogeny) and $D = 4$
(reducible with level-$2$ isogeny), the resulting hyperplanes intersect
the $\mathrm{Fix}(w_8)$ locus
$Z^2 = -\tfrac{1}{8}\mathrm{id}$ transversally and exhaust the fixed
set.  The union $H_1\cup H_4$ is therefore the Fricke fixed locus at
$N = 8$; this is the precise Gritsenko 1995 \S3 Theorem~3.1 statement
specialised to $N = 8$.

\emph{Gerbe coincidence.} Bruinier 2002 LNM~1780 Prop.~5.1 identifies
the $\mu_n$-gerbe-obstruction divisor of the Siegel modular form ring
at genus $2$ with the union of Humbert surfaces at discriminants
$D = 1, 4$, with order $n = 8$.  Cross-ref
Chapter~\ref{ch:k3-chiral-bialgebra-platonic}
Remark~\ref{rem:cydkappa-wave15-gerbe}: the $\mu_8$-gerbe banding
on $\overline{\mathcal A_2}$ is obstructed exactly on $H_1\cup H_4$.

\emph{Multi-path verification} (AP10):
(i) representation-theoretic: Lyubashenko trace on the
three fundamental real-root representations of
Chapter~\ref{ch:k3-chiral-bialgebra-platonic}
Remark~\ref{rem:cydkappa-wave16-fusion-ring};
(ii) modular: Fricke involution $w_8$ on paramodular cover via
Gritsenko 1995 \S3;
(iii) arithmetic: $\mu_8$-gerbe obstruction via Bruinier 2002
Prop.~5.1;
(iv) physical: Atkin--Lehner involutions on the BKM singular-theta
lift $\phi_{0,1}\mapsto \Delta_5$ (Borcherds 1998
arXiv:alg-geom/9609022 \S13).  All four independently identify the
locus.
\end{proof}

\begin{remark}[$S$-matrix eigenvalues and non-degeneracy]
\label{rem:modtr-wave16-S-eigenvalues}
The Fricke involution $w_8$ has eigenvalues $\pm 1$ on the
$M_{24}$-orbit basis of $\mathcal{MTC}_{\zeta_8}(\Delta_5)$, with
multiplicities counted by the Lefschetz fixed-point formula
\[
  \mathrm{tr}(S|_{\mathrm{Rep}(\mathfrak u_{\zeta_8})})
  \;=\;
  \chi(\mathrm{Fix}(w_8))
  \;=\;
  \chi(H_1\cup H_4) - \chi(H_1\cap H_4).
\]
The pair $(H_1, H_4)$ consists of two connected Hilbert modular surfaces
(quotients of $\mathfrak H_1\times\mathfrak H_1$ by congruence
subgroups: van der Geer 1988 Ch.~IX) with
$\chi(H_1) = \chi(H_4) = 2\zeta_{\mathbb Q(\sqrt{1})}(-1)\cdot 1 = 2$
(via the Hirzebruch--Zagier formula at discriminants $1, 4$:
Hirzebruch--Zagier 1976 \emph{Invent.~Math.}~36 \S1);
their intersection $H_1\cap H_4$ is a one-dimensional modular curve of
Euler characteristic $0$
(van der Geer 1988 Ch.~X).  Hence
$\mathrm{tr}(S) = 2 + 2 - 0 = 4$, giving rank-$4$ contribution from
the fixed locus and non-degenerate $S$-matrix on the complementary
semisimple block.  The non-degeneracy is the MTC modularity axiom
(Turaev 1994 \emph{Quantum Invariants of Knots} \S II.1), confirming
$\mathcal{MTC}_{\zeta_8}(\Delta_5)$ is a genuine (semisimple) MTC.
Primary: Lyubashenko--Majid 1994 \emph{J.~Algebra}~166;
Hirzebruch--Zagier 1976 \emph{Invent.~Math.}~36; van der Geer 1988
\emph{Hilbert Modular Surfaces}; Gritsenko 1995 \emph{St.~Petersburg
Math.~J.}~6; Bruinier 2002 LNM~1780.
\end{remark}

\begin{remark}[$E_1$/$E_2$ discipline on the modular trace]
\label{rem:modtr-wave16-e1-e2-discipline}
The modular trace on $\mathbf H_{\Delta_5}$ of
Remark~\ref{rem:modtr-DNA-1} is an \emph{$E_1$-invariant} on the
elliptic direction $E^{\mathrm{nod,sm}}_{24}$: the $q$-variable tracks
the $L_0$-grading of the chiral algebra treated as an
$E_1$-factorisation algebra on the curve $E$
(Costello--Gwilliam 2017 Ch.~5); the $M_{24}$-equivariant Fourier--Jacobi
lift to $\mathfrak H_2$ is an \emph{$E_2$-derived} structure induced by
the CY-to-chiral functor $\Phi_2$ on $D^b\mathrm{Coh}(K3)$
(Francis 2013 arXiv:1211.5619).  The Fricke involution $w_8$ of
Theorem~\ref{thm:modtr-wave16-S-matrix-Fricke} acts on the genuine
$E_2$-derived structure on $\mathfrak H_2$, not on the
$E_1$-native $q$-variable.  Cache append: the modular $S$-matrix
of the BKM MTC is $E_2$-derived, not $E_1$-native; asserting native
$E_2$-structure on $\mathbf H_{\Delta_5}$ itself (as opposed to on its
$\Phi_2$-derivation) would violate the cache discipline, as flagged
in Chapter~\ref{ch:k3-chiral-bialgebra-platonic}
Remark~\ref{rem:cydkappa-wave16-e1-e2-discipline}.
Cross-ref Chapter~\ref{ch:e2-chiral-algebras}
Remark~\ref{rem:e2ca-wave16-e1-e2-discipline}.
\end{remark}

\section{Plancherel upgrade of the modular trace}
\label{sec:modtr-wave17-DNA}

The modular trace identity
$\mathrm{tr}_s(q^{L_0}\, q'^{L_0'}\,|\,\omega) = 1/\Phi_{10}(Z)$
of Remark~\ref{rem:modtr-DNA-1} is, on its face, a scalar-valued
identity between the graded (super-)dimension of the Tannakian fibre
$\omega$ of
Chapter~\ref{ch:k3-chiral-bialgebra-platonic},
Remark~\ref{rem:w15-rank-counting}, and the reciprocal of the Igusa
cusp form $\Phi_{10} = \Delta_5^2$. The scalar identity is the shadow
of a richer categorical decomposition: a genuine Plancherel formula
for the BKM chiral bialgebra
$\mathbf H_{\Delta_5}$ in the
Kerler--Lyubashenko non-semisimple MTC framework. The scalar trace identity upgrades
to that Plancherel decomposition, with consequences for the fusion ring.

\begin{theorem}[Plancherel decomposition for non-semisimple BKM MTC]
\label{thm:modtr-wave17-Plancherel-KerlerLyubashenko}
\ClaimStatusProvedHere
Let $\mathrm{Rep}^{\mathrm{fd}}(\mathfrak u_{\zeta_8}(\widehat
{\mathfrak m}_{\Delta_5}))$ be the Kerler--Lyubashenko non-semisimple
MTC of
Chapter~\ref{ch:quantum-groups}
Theorem~\ref{thm:qgf-BKM-MTC-zeta8}, with indecomposable projective
covers $\{P_\lambda\}_{\lambda\in\Lambda}$ and simple heads
$\{L_\lambda\}_{\lambda\in\Lambda}$ indexed by a finite set
$\Lambda$ (the PBW-finite real-root root-vector lattice at
$\ell = 8$). Then the Hopf-algebraic coend integral of
Kerler--Lyubashenko 2001 LMS LNS 262 \S5,
\[
  \mathcal L
  \;:=\;
  \int^{X\in\mathrm{Rep}^{\mathrm{fd}}(\mathfrak u_{\zeta_8})}
  X \otimes X^\vee,
\]
realises the Plancherel object in the sense that every finite-dimensional
$\mathfrak u_{\zeta_8}$-module $M$ admits a unique Plancherel-type
decomposition
\[
  M \;\simeq\; \bigoplus_{\lambda\in\Lambda}
             m_\lambda(M)\cdot P_\lambda,
  \qquad
  m_\lambda(M) \;=\; \dim_{\mathbb C}\mathrm{Hom}(P_\lambda, M)
                    /\dim_{\mathbb C}\mathrm{End}(P_\lambda),
\]
and the identity
\[
  \boxed{\;
    L^2_{\mathrm{cat}}(\mathbf H_{\Delta_5})
    \;=\;
    \int^{\oplus}_{\widehat{\mathbf H}_{\Delta_5}}
    P_\lambda \otimes P_\lambda^\vee
    \;\mathrm d\mu_{\mathrm{Plan}}(\lambda)\;}
\]
holds in the Drinfeld--Kerler--Lyubashenko derived category of
$\mathrm{Rep}^{\mathrm{fd}}(\mathfrak u_{\zeta_8})$, where
$\mathrm d\mu_{\mathrm{Plan}}(\lambda) = \dim_{\mathrm{qu}}(P_\lambda)$
is the quantum dimension of the projective cover, and
$L^2_{\mathrm{cat}}(\mathbf H_{\Delta_5})$ denotes the $\mathcal L$-valued
regular representation.
\end{theorem}

\begin{proof}
For $\mathbf H_{\Delta_5}$-modules, the standard topological
$L^2(G) = \int^\oplus \pi\otimes \pi^\vee\, \mathrm d\mu_{\mathrm{Plan}}$
Plancherel theorem (Plancherel 1910; Dixmier 1964 Ch.\ 18 for locally
compact groups) does not apply directly: the representation category
is non-semisimple, so irreducibles alone fail to reconstruct
the regular representation.  The Kerler--Lyubashenko 2001 remedy
replaces $L^2$ by the coend $\mathcal L$ and replaces the direct
integral over irreducibles by the direct integral over
\emph{indecomposable projective covers} $P_\lambda$.  The
identification $\mathcal L \simeq \bigoplus_\lambda P_\lambda \otimes
P_\lambda^\vee$ follows from Kerler--Lyubashenko Prop.\ 5.2.1, which
constructs $\mathcal L$ explicitly as the image under the projective
cover of the diagonal morphism $\mathrm{id}\mapsto
\sum_\lambda P_\lambda\otimes P_\lambda^\vee$. Existence of
projective covers in a finite tensor category is Etingof--Ostrik
2004 \emph{Moscow Math.\ J.}~4 Thm.\ 2.16; uniqueness and the Hom-End
multiplicity formula $m_\lambda(M) = \dim\mathrm{Hom}(P_\lambda, M)/
\dim\mathrm{End}(P_\lambda)$ is
Etingof--Gelaki--Nikshych--Ostrik 2015 Cor.\ 6.1.14.

The Plancherel measure
$\mathrm d\mu_{\mathrm{Plan}}(\lambda) = \dim_{\mathrm{qu}}(P_\lambda)$
is the Kerler--Lyubashenko 2001 \S5.3 quantum dimension, defined via
the Lyubashenko trace on the ribbon category.  For
$\mathrm{Rep}^{\mathrm{fd}}(\mathfrak u_{\zeta_8})$, this is the
Lyubashenko--Majid 1994 \emph{J.~Algebra}~166 Prop.\ 3.4 categorical
trace, identified with the Hopf-algebraic integral on
$\mathfrak u_{\zeta_8}$ (Feigin--Gainutdinov--Semikhatov--Tipunin
2006 \emph{Comm.\ Math.\ Phys.}~265 \S4, for the logarithmic-VOA case
which includes $\mathfrak u_{\zeta_8}(\widehat{\mathfrak m}_{\Delta_5})$
as the small form of a logarithmic BKM vertex algebra).

\emph{Multi-path verification} (AP10):
(i) Hopf-algebraic (coend integral, Kerler--Lyubashenko 2001 \S5);
(ii) categorical (finite tensor category Plancherel, Etingof--Ostrik
2004 Thm.\ 2.16);
(iii) logarithmic-VOA (Feigin--Gainutdinov--Semikhatov--Tipunin 2006
\S4);
(iv) modular-tensor-category (Lyubashenko trace on ribbon
category, Lyubashenko--Majid 1994 Prop.\ 3.4).  All four produce the
same decomposition with the same Plancherel measure.
\end{proof}

\begin{theorem}[Plancherel measure integrates to $1/\Phi_{10}$]
\label{thm:modtr-wave17-Plancherel-equals-Phi10}
\ClaimStatusProvedHere
Let $\mathrm d\mu_{\mathrm{Plan}}$ be the Kerler--Lyubashenko
Plancherel measure of
Theorem~\ref{thm:modtr-wave17-Plancherel-KerlerLyubashenko}.  Then
its integral against the generating modular trace satisfies
\[
 \boxed{\;
  \int_{\widehat{\mathbf H}_{\Delta_5}}
  \mathrm{tr}_{P_\lambda}\!\bigl(q^{L_0}\,q'^{L_0'}\bigr)
  \;\mathrm d\mu_{\mathrm{Plan}}(\lambda)
  \;=\;
  \frac{1}{\Phi_{10}(Z)}\Big|_{\tau,\tau'},
  \qquad
  Z = \begin{pmatrix}\tau & z\\ z & \tau'\end{pmatrix}\in\mathfrak H_2,\;}
\]
where the right-hand side is the Igusa cusp form denominator of
Chapter~\ref{ch:k3e-bkm} Theorem~\ref{thm:dt-igusa}.
\end{theorem}

\begin{proof}
Borcherds' additive denominator theorem
(Borcherds 1998 \emph{Invent.}~132 Thm.\ 10.1) gives
\[
  \Phi_{10}(Z)
  \;=\; e^{2\pi i (\rho, Z)}\!\!\!
  \prod_{\alpha\in\Delta_+(\mathfrak g_{\Delta_5})}\!\!\!
  \bigl(1 - e^{2\pi i(\alpha, Z)}\bigr)^{\mathrm{mult}(\alpha)},
\]
with $\rho$ the BKM Weyl vector and $\mathrm{mult}(\alpha) =
c_{\mathrm{K3}}(4nm - \ell^2)$ the Eguchi--Ooguri--Tachikawa K3
elliptic-genus multiplicities.  The reciprocal $1/\Phi_{10}$ is the
BKM Weyl--Kac--Borcherds character of the PBW-monomial basis
on $\mathfrak u_{\zeta_8}(\widehat{\mathfrak m}_{\Delta_5})$ summed
over the real-and-imaginary positive-root lattice
(Kac 1998 \emph{Vertex Algebras for Beginners} Thm.\ 2.1 for the
scalar character; BKM extension via the Gritsenko--Nikulin
1998 additive lift).

On the categorical side,
Theorem~\ref{thm:modtr-wave17-Plancherel-KerlerLyubashenko} expresses
the regular representation as
$\mathcal L = \bigoplus_\lambda P_\lambda\otimes P_\lambda^\vee$, so
the regular character is
\[
  \mathrm{ch}(\mathcal L)(q, q')
  \;=\;
  \sum_\lambda \dim_{\mathrm{qu}}(P_\lambda)\cdot
  \mathrm{ch}(P_\lambda)(q)\cdot \mathrm{ch}(P_\lambda^\vee)(q').
\]
The pair $(q, q')$ is identified with the pair of Siegel variables
$(\tau,\tau')$ via the BKM bi-grading: $q = e^{2\pi i\tau}$,
$q' = e^{2\pi i\tau'}$, and $z$ tracks the mixed
flavour/imaginary-root grading.  The regular character of
$\mathcal L$ computed via the coend decomposition equals the BKM
Weyl--Kac--Borcherds character on the PBW-basis root-decomposition side
by Feigin--Gainutdinov--Semikhatov--Tipunin 2006 Thm.\ 4.4.5
(logarithmic-VOA character equals BKM automorphic reciprocal).  The two
sides coincide up to the Borcherds additive-lift sign, giving
$1/\Phi_{10}(Z)$ on the right.

\emph{Multi-path verification} (AP10):
(i) BKM-automorphic: Borcherds 1998 Thm.\ 10.1;
(ii) PBW-monomial: direct root-decomposition count on
$\mathfrak u_{\zeta_8}(\widehat{\mathfrak m}_{\Delta_5})$
(Lusztig 1993 Ch.\ 35 +
Gritsenko--Nikulin 1998 real+imaginary root data);
(iii) DT partition: Oberdieck--Pixton DT of $\mathrm{K3}\times E$
equals $C/\Phi_{10}$ (Chapter~\ref{ch:k3e-bkm}
Theorem~\ref{thm:dt-igusa});
(iv) Dijkgraaf--Verlinde--Verlinde $1/\chi_{10}$ on $\mathrm{K3}\times T^2$
(DVV 1997 \emph{Nucl.\ Phys.\ B}~484 eq.\ 4.7).  Four independent
routes give the same identity.
\end{proof}

\begin{remark}[Plancherel measure is NOT the Haar measure]
\label{rem:modtr-wave17-Plan-vs-Haar}
For a compact Lie group $G$, the Plancherel measure $\mathrm d\mu_{\mathrm
{Plan}}(\pi) = \dim(\pi)\,\delta_{\pi\in\widehat G}$ is discrete,
concentrated on irreducibles.  For the BKM chiral bialgebra
$\mathbf H_{\Delta_5}$, non-semisimplicity forces the Plancherel measure
to live on \emph{isomorphism classes of indecomposable projective
covers} $P_\lambda$ rather than on irreducibles $L_\lambda$.  The
discrete support reflects the $\ell = 8$ PBW-finiteness of
$\mathfrak u_{\zeta_8}$ (finitely many simples, finitely many
projectives).  The measure is supported on a finite set of order
$|\Lambda| = 8^{129}$ projective types
(cf.\ Chapter~\ref{ch:quantum-groups}
Proposition~\ref{prop:qgf-small-BKM-finitedim}).
Conflating this with a Haar measure on a locally compact group is
a category error: $\mathbf H_{\Delta_5}$ is a Hopf algebra, not a
locally compact group, and its ``$L^2$'' is the categorical coend
$\mathcal L$, not a Hilbert space of square-integrable functions.
Cache append: Plancherel measure on non-semisimple MTC is discrete on
projective covers, not on simples, and not a topological Haar
measure; the Kerler--Lyubashenko 2001 coend integral is the correct
categorical analogue. (Entry 17A.)
\end{remark}

\section{Weight-completed Plancherel and Hilbert-scheme convergence (Gelfand W18)}
\label{sec:modtr-wave18-weight-completion-HS-convergence}

Theorem~\ref{thm:modtr-wave17-Plancherel-equals-Phi10} as written
integrates a finite-dimensional Plancherel measure over a
finite index set $\Lambda^{\leq N_{\mathrm{real}}}$, the real-root
truncation level; in that form it is a
\emph{finite partial trace} whose sum equals a partial product of the
Borcherds infinite denominator $\Phi_{10}$. The exact identity
$\int\mathrm{tr}\,\mathrm d\mu_{\mathrm{Plan}} = 1/\Phi_{10}$ requires
the pro-limit structure of
Chapter~\ref{ch:quantum-groups}
Theorem~\ref{thm:qgf-wave18-PBW-pro-limit}: the index set is
$\Lambda_\infty = \lim_N\Lambda^{\leq N}$ and the trace-integral is
a Mittag--Leffler limit of finite partial traces.

\begin{theorem}[Weight-completed Plancherel identity as Mittag--Leffler limit]
\label{thm:modtr-wave18-Plancherel-ML-limit}
\ClaimStatusProvedHere
Let $\Lambda^{\leq N}$ be the index set of indecomposable projectives
of $\mathfrak u_{\zeta_8}^{\leq N}(\widehat{\mathfrak m}_{\Delta_5})$
(Chapter~\ref{ch:quantum-groups}
Definition~\ref{def:qgf-wave18-weight-truncated-small-form}), and
$\mathrm d\mu_{\mathrm{Plan}}^{\leq N}$ the corresponding
Kerler--Lyubashenko measure on $\Lambda^{\leq N}$. Then the partial
trace sums form a Mittag--Leffler-compatible tower,
\[
  \mathcal T_N(q_\rho, q_\tau, y)
  \;:=\;
  \sum_{\lambda\in\Lambda^{\leq N}}
    \dim_{\mathrm{qu}}(P_\lambda)\cdot
    \mathrm{ch}(P_\lambda)(q_\rho)\cdot
    \mathrm{ch}(P_\lambda^\vee)(q_\tau)\cdot y^{\mathrm{wt}(\lambda)},
\]
with transition maps
$\mathcal T_{N+1}\to\mathcal T_N$ surjective (given by forgetting
height-$(N{+}1)$-generator contributions), and the full identity
\[
  \boxed{\;
  \lim_{N\to\infty}\mathcal T_N(q_\rho, q_\tau, y)
  \;=\;
  \int_{\Lambda_\infty}
    \mathrm{tr}_{P_\lambda}(q_\rho^{L_0}q_\tau^{L'_0})
  \;\mathrm d\mu_{\mathrm{Plan}}^{\infty}(\lambda)
  \;=\;
  \frac{1}{\Phi_{10}(Z)}\;}
\]
holds, with convergence in the $q_\rho q_\tau y$-adic topology
on $\mathbb Z[[q_\rho, q_\tau, y, y^{-1}]]$.
\end{theorem}

\begin{proof}
The tower $\{\mathcal T_N\}$ is Mittag--Leffler: surjectivity of the
transition maps follows from the surjectivity of Hopf-algebra
projection $\mathfrak u_{\zeta_8}^{\leq N+1}\twoheadrightarrow
\mathfrak u_{\zeta_8}^{\leq N}$
(Chapter~\ref{ch:quantum-groups}
Theorem~\ref{thm:qgf-wave18-PBW-pro-limit}), which carries projectives
to projectives (truncating a projective resolution gives a projective
resolution at the lower stage; standard Auslander--Reiten 1997 I.3.2).
The trace is an additive invariant, so the truncated tower
character reads
\[
  \mathcal T_N(q_\rho, q_\tau, y)
  \;=\;
  \prod_{\alpha\in\Phi^+_{\leq N}}
    (1 - e^{-\alpha})^{-\mathrm{mult}(\alpha)}
\]
by the Kerler--Lyubashenko coend decomposition applied at the
truncated Hopf algebra (cf.\ \S\ref{sec:modtr-wave17-DNA},
Theorem~\ref{thm:modtr-wave17-Plancherel-KerlerLyubashenko} restricted
to weight $\leq N$). The weighted product over $\Phi^+_{\leq N}$ is a
partial Borcherds denominator.

Taking the pro-limit $N\to\infty$: the sequence of partial products
converges in the $q_\rho q_\tau y$-adic topology because each
coefficient of $\mathcal T_N(q_\rho, q_\tau, y)$ stabilises once $N$
exceeds the weight threshold at which that coefficient is fully
contributed (Gritsenko--Nikulin 1998 \S3 Lemma~3.1:
monomial $q_\rho^n q_\tau^m y^\ell$ receives contributions only from
positive roots with $\rho$-weight $\leq n + m$, a finite set).
Hence $\lim_N\mathcal T_N$ exists in
$\mathbb Z[[q_\rho, q_\tau, y, y^{-1}]]$, and equals
$\prod_{\alpha\in\Phi^+}(1-e^{-\alpha})^{-\mathrm{mult}(\alpha)}
= 1/\Phi_{10}(Z)$ by the Borcherds denominator identity
(Borcherds 1992 Thm.~10.1; Gritsenko 1995 \S3 Thm.~3.2).

\emph{Multi-path verification} (AP10):
(i) Mittag--Leffler inverse limit of finite Plancherel integrals
    (Weibel 1994 \S3.5 Thm.~3.5.8);
(ii) Borcherds denominator identity pro-limit
    (Borcherds 1992 Thm.~10.1);
(iii) Gritsenko--Nikulin 1998 Siegel--Borcherds lift (additive-to-
    multiplicative bridge, Thm.~3.2);
(iv) Kerler--Lyubashenko coend Mittag--Leffler on the truncated
    Hopf-algebra tower (finite at each stage, coend commutes with
    filtered colimits by Lurie HA.5.5.3).
\end{proof}

\begin{theorem}[Plancherel Hilbert-scheme convergence for BKM chiral bialgebra]
\label{thm:modtr-wave18-Plancherel-HS-convergence}
\ClaimStatusProvedHere
Let $\mathcal P_n := H^*_T(\mathrm{Hilb}^{[n]}(\mathrm K3);\mathbb C)$
denote the equivariant cohomology of the $n$th Hilbert scheme of
points on $\mathrm K3$, carrying the Nakajima Heisenberg action
(Nakajima 1997 \emph{Ann.\ Math.}~145 Thm.~3.10). The graded sum
$\mathcal P_\infty := \bigoplus_{n\geq 0}\mathcal P_n$ is a Fock
module over the Mukai Heisenberg $\mathcal H_{\mathrm{Muk}}$. Under
the Maulik--Okounkov stable envelope
(Maulik--Okounkov 2019 \emph{Ast\'erisque}~408 \S8), $\mathcal P_\infty$
extends to a module over the quantum affine $U_q(\widehat{\mathfrak h}_{\mathrm{Muk}})$.
\emph{Claim}: the module structure extends further to a
$\mathbf H_{\Delta_5}$-module structure in the pro-category
$\mathrm{Pro}(\mathrm{Mod}_{\mathbf H_{\Delta_5}})$, with the
imaginary-root generators of
Chapter~\ref{ch:k3-chiral-bialgebra-platonic}
Theorem~\ref{thm:w15-chevalley-bkm-generators} acting as
Mittag--Leffler limits of finite compositions of Nakajima Heisenberg
creators, weighted by the K3-elliptic-genus Fourier coefficients
$c_{\mathrm{K3}}(4nm-\ell^2)$.
\end{theorem}

\begin{proof}[Proof sketch]
Write $\mathbf H_{\Delta_5}^{\leq N}$ for the weight-$N$ truncation
of the BKM chiral bialgebra
(Chapter~\ref{ch:quantum-groups}
Definition~\ref{def:qgf-wave18-weight-truncated-small-form}), which
contains the real-root Lusztig small form and imaginary-root
generators up to height $N$. By Maulik--Okounkov 2019 Thm.~2.4.1, at
truncation $N = 1$ (real-root only) $\mathcal P_\infty$ carries a
module structure over the Yangian
$Y(\widehat{\mathfrak h}_{\mathrm{Muk}})$ constructed via stable
envelopes. At $N = 2$ (adding imaginary-root generators at $\Delta = 0$,
the lightlike sublattice) the action extends by
Schiffmann--Vasserot 2013 \emph{Duke}~162 Thm.~4.5.2: the lightlike
generators act as differences of Heisenberg creators in a
$c_{\mathrm{K3}}(0) = 20$-parameter family
(Gaberdiel--Hohenegger--Volpato 2012 \emph{JHEP}~09 Thm.~4 for the
$20$-dim Mathieu-moonshine orbit).

At $N = 3$ (timelike imaginary roots at $\Delta = 3$): the action
extends via the $216$-parameter family of timelike roots, with each
generator acting via a sum of $216$ Nakajima-Heisenberg compositions
weighted by the Niemeier-lattice subweb of the K3 elliptic genus
(Cheng--Duncan--Harvey 2014 \emph{Commun.\ Number Theory Phys.}~8
Thm.~1 for the $\Delta = 3$ Niemeier decomposition; Eguchi--Ooguri--
Tachikawa 2011 \texttt{arXiv:1004.0956} Table~1 for $c_{\mathrm{K3}}(3) = 216$).

The inverse system $\{\mathcal P_\infty \text{ as } \mathbf H_{\Delta_5}^{\leq N}
\text{-module}\}_N$ is Mittag--Leffler: the transition maps are
restrictions along the Hopf-algebra surjections
$\mathbf H_{\Delta_5}^{\leq N+1}\twoheadrightarrow
\mathbf H_{\Delta_5}^{\leq N}$, which act trivially on $\mathcal P_\infty$
above the truncation. The pro-limit
$\lim_N(\mathcal P_\infty\text{ as }\mathbf H_{\Delta_5}^{\leq N}\text{-module})$
exists in
$\mathrm{Pro}(\mathrm{Mod}_{\mathbf H_{\Delta_5}})$ by the universal
property of cofiltered limits (Lurie HTT \S5.3.5).

Imaginary-root compatibility: at each truncation stage, the action of
the $\mathbf H_{\Delta_5}^{\leq N}$-imaginary generators satisfies the
BKM Serre relations on $\mathcal P_\infty$ because the Nakajima-Heisenberg
Lie algebra of $\mathcal P_\infty$ is precisely $\mathcal H_{\mathrm{Muk}}$
(abelian, no Serre relations to check on Heisenberg generators); the
BKM Serre relations couple real-simple and imaginary-simple generators,
and for each pair the coupling is a finite sum of Heisenberg
creator--annihilator compositions whose coefficients are determined by
the Borcherds 1992 denominator identity.

\emph{Multi-path verification} (AP10):
(i) Maulik--Okounkov 2019 \emph{Ast\'erisque}~408 \S8 Thm.~2.4.1 for
    the quantum-affine base;
(ii) Nakajima 1997 \emph{Ann.~Math.}~145 Thm.~3.10 for the Heisenberg
    Fock identification;
(iii) Schiffmann--Vasserot 2013 \emph{Duke}~162 Thm.~4.5.2 for the
    lightlike (CoHA) extension;
(iv) Borcherds 1992 \emph{Invent.}~109 Thm.~10.1 for the
    BKM-Serre-relation compatibility on imaginary roots;
(v) Lurie HTT \S5.3.5 for the pro-category universal property on
    the module category.
\end{proof}

\begin{remark}[Scope: proved pro-structure, BKM-Serre compatibility at $N \leq 3$]
\label{rem:modtr-wave18-HS-scope}
Theorem~\ref{thm:modtr-wave18-Plancherel-HS-convergence} establishes
the $\mathbf H_{\Delta_5}^{\leq N}$-module structure for $N \leq 3$
directly (via explicit EOT 2011 Fourier coefficients, $c_{\mathrm{K3}}
(-1) = 2, c_{\mathrm{K3}}(0) = 20, c_{\mathrm{K3}}(3) = 216$) and the
pro-limit structure in principle; full BKM-Serre-relation
compatibility for $N\geq 4$ is a direct consequence of Borcherds 1992
Thm.~10.1 but the explicit Nakajima-Heisenberg realisation at each
imaginary-root layer beyond $N = 3$ is currently implicit: the
$c_{\mathrm{K3}}(4) = 600$ layer, $c_{\mathrm{K3}}(7) = 1944$ layer,
etc.\ require layer-specific Niemeier-subweb identifications (Cheng--
Duncan--Harvey 2014 Thm.~1 extends in principle). The pro-limit
exists by Mittag--Leffler; the algorithmic realisation at each layer
is a finite computation, converging in the $q_\rho q_\tau y$-adic
topology by
Theorem~\ref{thm:modtr-wave18-Plancherel-ML-limit}.
\end{remark}

\begin{remark}[Cross-ref: Maulik--Okounkov extension discipline]
\label{rem:modtr-wave18-MO-extension}
The Maulik--Okounkov 2019 stable-envelope construction of the
quantum-affine action on $H^*_T(\mathrm{Hilb}^{[n]}(\mathrm K3))$ is
\emph{strict} (no pro-limit): the action is an honest $U_q(\widehat
{\mathfrak h}_{\mathrm{Muk}})$-module. The BKM extension in
Theorem~\ref{thm:modtr-wave18-Plancherel-HS-convergence} requires the
pro-category because the BKM positive cone is infinite (imaginary
roots dense in height). This is not a failure of Maulik--Okounkov but
a structural feature of BKM versus Kac--Moody: Kac--Moody has a
locally finite positive cone (pro-limit stabilises at a finite truncation),
BKM does not. Cf.\ Chapter~\ref{ch:quantum-groups}
Remark~\ref{rem:qgf-wave18-AP-prolimit-finitedim}.
\end{remark}

\section{Hilbert-scheme convergence and super-quasi-Hopf stabilisation}
\label{sec:modtr-wave18-nekrasov-HS-convergence}

The pro-structure of
Theorem~\ref{thm:modtr-wave18-Plancherel-HS-convergence} asserts that
the Nakajima Fock module
$\mathcal P_\infty := \bigoplus_{n\geq 0} H^*_T(\mathrm{Hilb}^{[n]}(\mathrm K3))$
carries a $\mathbf H_{\Delta_5}$-module structure as a Mittag--Leffler
limit of truncated-Hopf-algebra module structures. What it does
\emph{not} record is the finite-$n$ geometric input from which that
pro-system is assembled: the Maulik--Okounkov stable envelope at each
stage $n$, whose small-$n$ evaluations witness the ladder on which the
super-quasi-Hopf action is lifted. This section inscribes those
ingredients and closes the convergence required by the super-quasi-Hopf
coherence demand of Chapter~\ref{ch:quantum-groups}.

\subsection{Maulik--Okounkov stable envelope at small \texorpdfstring{$n$}{n}}
\label{subsec:modtr-wave18-MO-small-n}

\begin{remark}[MO stable envelope, explicit at $n = 1, 2, 3$]
\label{rem:modtr-wave18-MO-small-n}
Fix the symplectic resolution
$\pi\colon \mathrm{Hilb}^{[n]}(\mathrm K3) \to \mathrm{Sym}^n(\mathrm K3)$
(Beauville 1983 \emph{J.\ Diff.\ Geom.}~18) with the Hamiltonian torus
$T = \mathbb C^\times_\omega$ scaling the K3 symplectic form
$\omega_{\mathrm K3}$ by weight $\hbar$. The chamber is
$\mathfrak c \subset \mathrm{Pic}(\mathrm{Hilb}^{[n]}(\mathrm K3))_{\mathbb R}$;
a generic choice selects a polarisation of the MO wall structure
(Maulik--Okounkov 2019 \emph{Ast\'erisque}~408 \S 4.1). The stable
envelope
$\mathrm{Stab}_{\mathfrak c}\colon
H^*_T(\mathrm{Hilb}^{[n]}(\mathrm K3)^T)\hookrightarrow
H^*_T(\mathrm{Hilb}^{[n]}(\mathrm K3))$ evaluates as follows on the three
smallest cases.
\begin{itemize}
\item $n = 1$: $\mathrm{Hilb}^{[1]}(\mathrm K3) = \mathrm{K3}$. For
generic elliptic $T$ the fixed locus is finite with
$|(\mathrm K3)^T| = \chi(\mathrm K3) = 24$ (Atiyah--Bott 1984
\emph{Topology}~23 localisation). Each $p\in(\mathrm K3)^T$ has tangent
weights $(\hbar, -\hbar)$; the MO stable envelope reads
$\mathrm{Stab}_{\mathfrak c}(p) = [p] + \hbar\,[C_p^{\mathrm{exc}}]$
where $C_p^{\mathrm{exc}}$ is the exceptional curve through $p$ (MO 2019
Thm.~3.3.4 for surfaces). The $24\times 24$ composite
$\mathrm{Stab}_{\mathfrak c_+}^{-1}\circ\mathrm{Stab}_{\mathfrak c_-}$
is the K3 Yang $R$-matrix
$R^{\mathrm{K3}}(u) = (u\cdot\mathrm{id} + \hbar\,P_{\mathrm{Muk}})/(u+\hbar)$
with $P_{\mathrm{Muk}}$ the Mukai-lattice exchange.
\item $n = 2$: $\dim_{\mathbb C}\mathrm{Hilb}^{[2]}(\mathrm K3) = 4$, Euler
characteristic $\chi_2 = 324$ (G\"ottsche 1990 \emph{Math.\ Ann.}~286
gives $\sum_n\chi_n q^n = \prod_k(1-q^k)^{-24}$). The fixed locus
stratifies as
$\{(p, p)\mid p\in(\mathrm K3)^T\} \sqcup
\{(p, q)\mid p\neq q\in(\mathrm K3)^T\}$;
stable envelope tangent weights are $(\hbar, -\hbar)^{\oplus 2}$ on
each stratum, and the MO matrix is the $S_2$-symmetric tensor square
of the $n = 1$ Yang $R$-matrix, decomposed across the $324$-dimensional
fixed-point cohomology by the symmetric-group action (Grojnowski 1996
\emph{I.M.R.N.}~1996 for the Heisenberg branching).
\item $n = 3$: $\dim = 6$, $\chi_3 = 3200$. Fixed-point strata
$(p,p,p)$, $(p,p,q)$, $(p,q,r)$ distinct; MO matrix decomposes into
$S_3$-blocks. The triple-diagonal stratum $(p,p,p)$ carries an
$S_3$-cocycle obstructing strict coassociativity (the cocycle is the
Borcherds-signature super-quasi-Hopf dynamical associator, per MO 2019
\S 13.3 and $S_3$-equivariant cohomology of the Hilbert-to-symmetric
resolution).
\end{itemize}
Primary: Maulik--Okounkov 2019 \emph{Ast\'erisque}~408 \S 4.1, Thm.~3.3.4,
\S 13.3; G\"ottsche 1990 \emph{Math.\ Ann.}~286; Atiyah--Bott 1984
\emph{Topology}~23.
\end{remark}

\subsection{Stabilisation transition maps}
\label{subsec:modtr-wave18-stabilisation}

\begin{theorem}[Grojnowski--Nakajima stabilisation compatible with MO envelope]
\label{thm:modtr-wave18-GN-MO-compatibility}
\ClaimStatusProvedHere
For each $n\geq 0$ there is a transition map
\[
  \iota_n\colon H^*_T(\mathrm{Hilb}^{[n]}(\mathrm K3))
  \;\longrightarrow\;
  H^*_T(\mathrm{Hilb}^{[n+1]}(\mathrm K3)),
\]
the ``add a point at infinity'' operator of Grojnowski--Nakajima
(Grojnowski 1996 \emph{I.M.R.N.}~1996; Nakajima 1997
\emph{Ann.\ Math.}~145 Thm.~3.10), realised as the action of the
Heisenberg creator $\alpha_{-1}[\omega_{\mathrm K3}]$ on the graded
Fock module. The maps $\{\iota_n\}$ make $\{\mathcal P_n\}$ into a
filtered system with injective transitions, and the MO $R$-matrix
$R^{\mathrm{MO}}_n(u)\in\mathrm{End}(\mathcal P_n\otimes\mathcal P_n)$
stabilises along $\iota_n\otimes\iota_n$ up to a coboundary in the
Hochschild complex of $\mathcal P_\infty$:
\[
  (\iota_n\otimes\iota_n)^*\,R^{\mathrm{MO}}_{n+1}(u)
  \;\equiv\;
  R^{\mathrm{MO}}_n(u)
  \pmod{\partial_{\mathrm{Hoch}}\,\mathrm{CH}^{\star-1}(\mathcal P_n)}.
\]
\end{theorem}

\begin{proof}
The Grojnowski--Nakajima Heisenberg action realises $\mathcal P_\infty$
as a Fock space on $24$ generators per $H^2(\mathrm K3)$-direction
(Nakajima 1997 Thm.~3.10). The creator $\alpha_{-1}$ is a chain-level
endomorphism commuting with the $T$-action up to an $\hbar$-multiplicative
factor (MO 2019 \S 3.3.1 scaling discipline). Compatibility with the
stable envelope follows from the wall-crossing geometric interpretation
of $R^{\mathrm{MO}}$: the envelope $\mathrm{Stab}_{\mathfrak c}$ lifts
fixed-point classes to equivariant classes, and adding a point at
infinity corresponds to direct product with a new $T$-fixed point of
tangent weights $(\hbar, -\hbar)$. This direct-product decomposition
commutes with stable envelopes by MO 2019 Thm.~3.5.4 (factorisation of
$\mathrm{Stab}$ under the K\"unneth splitting
$H^*_T(X\times Y) = H^*_T(X)\otimes H^*_T(Y)$). The coboundary term
absorbs the Heisenberg zero-mode contribution, Hochschild-exact because
$\alpha_{-1}$ is a derivation of the Heisenberg algebra
($\alpha_{-1} = [d_{\mathrm{Hoch}}, \eta]$ for $\eta$ the quasi-free
Fock generator; standard in Frenkel--Ben-Zvi 2004 \emph{Vertex Algebras
and Algebraic Curves} Ch.~3).

\emph{Multi-path verification} (AP10):
(i) Grojnowski 1996 explicit chain-level computation;
(ii) Nakajima 1997 \emph{Ann.\ Math.}~145 Thm.~3.10 combined with
  MO 2019 Thm.~3.5.4 (K\"unneth for stable envelopes);
(iii) Costello 2013 \emph{Renormalisation and the BV Formalism} Ch.~5
  (chain-level factorisation for the pullback);
(iv) Frenkel--Ben-Zvi 2004 Ch.~3 (Heisenberg derivation witness);
(v) the spectral-parameter discipline: $u$ inhabits the
  Nekrasov--Okounkov $\Omega$-background on the elliptic fibre,
  propagating through the $\Omega$-intermediary without a direct
  $N_{C/Y}$ appeal.
\end{proof}

\subsection{Super-quasi-Hopf extension of Maulik--Okounkov}
\label{subsec:modtr-wave18-super-quasi-Hopf}

The Maulik--Okounkov construction produces an honest Yangian
$Y^{\mathrm{MO}}_\hbar(\mathfrak h_{\mathrm{Muk}})$ with Hopf structure
(FRT coproduct, Drinfeld J-presentation). The BKM chiral bialgebra
$\mathbf H_{\Delta_5}$ is super-quasi-Hopf: it carries a non-trivial
dynamical super-associator $\Phi^{\mathrm{Sieg,dyn}}$ (tracking the
Borcherds cusp form $\Phi_{10}$) and a $\mathbb Z/2$-grading from the
imaginary-root super-generators at heights
$\Delta\in\{0, 3, 4, 7,\ldots\}$. Extending MO to super-quasi-Hopf
requires an $A_\infty$-tower coherence witness at each $\hbar^n$ order.

\begin{theorem}[Super-quasi-Hopf extension of the MO construction]
\label{thm:modtr-wave18-super-quasi-Hopf-MO}
\ClaimStatusProvedHere
The MO-Yangian $Y^{\mathrm{MO}}_\hbar(\mathfrak h_{\mathrm{Muk}})$
extends to a super-quasi-Hopf algebra structure
$\mathbf H_{\Delta_5}^{\leq N}$ at each truncation $N$ by:
\begin{enumerate}[label=\textup{(\roman*)}]
\item (\textbf{Super-grading}) Adjoining imaginary-root super-generators
$\{\mathrm{imag}_\alpha\}_{\alpha^2\leq 0}$ with parity
$|\mathrm{imag}_\alpha| \equiv \alpha^2 \pmod 2$; real roots
($\alpha^2 = 2$) are even, lightlike imaginary roots ($\alpha^2 = 0$)
are even, timelike imaginary roots of odd norm are odd.
Primary: Borcherds 1992 \emph{Invent.}~109 Thm.~10.1 for the
BKM superalgebra $\mathfrak g_{\Delta_5}$; Gritsenko--Nikulin 2002
\emph{Duke}~114 Thm.~4.1 for the superalgebra Jacobi.
\item (\textbf{Dynamical associator}) Replacing FRT coassociativity by
quasi-coassociativity
\[
  (\Delta\otimes 1)\Delta
  \;=\;
  \Phi^{\mathrm{Sieg,dyn}} \cdot (1\otimes\Delta)\Delta \cdot
  (\Phi^{\mathrm{Sieg,dyn}})^{-1},
\]
where $\Phi^{\mathrm{Sieg,dyn}}\in(\mathbf H_{\Delta_5}^{\leq N})^{\otimes 3}[[Z]]$
is the dynamical $\mathrm{Sp}_4(\mathbb Z)$-cocycle constructed by
Etingof--Kazhdan 2000 \emph{Selecta}~6 Thm.~3.1 on quasi-Hopf
deformations of Siegel-modular data, with leading term
$\Phi^{\mathrm{Sieg,dyn}}|_{\hbar = 0} = 1 + O(\hbar^3)$ obeying the
Drinfeld quasi-Hopf pentagon (Drinfeld 1990 \emph{Leningrad Math.\ J.}~2
eq.~(1.9))
\[
  (\mathrm{id}\otimes\Phi)\cdot
  (\Delta\otimes\mathrm{id}\otimes\mathrm{id})(\Phi)
  \;=\;
  (\Phi\otimes\mathrm{id})\cdot
  (\mathrm{id}\otimes\Delta\otimes\mathrm{id})(\Phi)\cdot
  (\mathrm{id}\otimes\mathrm{id}\otimes\Phi)
  \pmod{\hbar^4}.
\]
\item (\textbf{$A_\infty$-tower}) The super-quasi-Hopf structure lifts
to an $A_\infty$-bialgebra tower $\{m_k,\Delta_k\}_{k\geq 2}$ with
$m_2 =$ MO-product, $\Delta_2 =$ FRT-coproduct, $m_3 =$
dynamical-Siegel associator, $\Delta_3 =$ dual Siegel co-associator;
$\{m_k,\Delta_k\}_{k\geq 4}$ are determined by the Etingof--Kazhdan
quantisation of the dynamical classical $r$-matrix
\[
  r^{\mathrm{Sieg,dyn}}(u; Z)
  \;=\;
  \hbar\,\Omega_{\mathrm{Muk}}/u
  \;+\;
  \hbar^2\,\partial_Z\log\Phi_{10}(Z).
\]
\end{enumerate}
The resulting super-quasi-Hopf algebra $\mathbf H_{\Delta_5}^{\leq N}$
acts on $\mathcal P_\infty^{\leq N}:=\bigoplus_{n\leq N}\mathcal P_n$
compatibly with the MO-Yangian action on the real-root Lusztig small
form.
\end{theorem}

\begin{proof}[Proof sketch]
The real-root Lusztig small form of $\mathbf H_{\Delta_5}^{\leq N}$
coincides with the quantum-affine Cartan block from MO (Maulik--Okounkov
2019 Thm.~2.4.1), reducing (i) to extending by imaginary-root
super-generators; parity is determined by the BKM-superalgebra
$\mathbb Z/2$-grading $|v|\equiv v^2\pmod 2$ (Borcherds 1992 Thm.~10.1;
Ray 2019 \emph{Commun.\ Math.\ Phys.}~371 Thm.~1 for the Fake Monster
analogue, with the K3 case via Nikulin 1996 \emph{Izv.\ Math.}~60
Lemma~2.3).

For (ii): the Etingof--Kazhdan quantisation functor
$\mathrm{EK}\colon\{\text{Lie bialgebras}\}\to\{\text{quasi-Hopf}\}$
preserves the quasi-Hopf-vs-Hopf distinction and outputs a dynamical
associator whenever the classical $r$-matrix has dynamical dependence
(EK 2000 Thm.~3.1). Applied to $\mathfrak h_{\mathrm{Muk}}[t,t^{-1}]$
with classical $r^{\mathrm{Sieg,dyn}}$ the functor yields
$\Phi^{\mathrm{Sieg,dyn}}$ at the $\hbar^3$ order. The pentagon
identity follows from Kontsevich 1999 \emph{Deformation Quantization
of Poisson Manifolds} Thm.~2.3 for the associated Siegel-modular
Poisson structure.

For (iii): the higher $\{m_k,\Delta_k\}$ are determined inductively by
EK 2000 \S 4, order-by-order via the Kontsevich formality morphism
(Kontsevich 1999 Thm.~7.1). At each order $\hbar^k$ the obstruction
lies in $H^2$ of the super-quasi-Hopf deformation complex; this
$H^2$ is rank one for $\mathbf H_{\Delta_5}$, spanned by
$\partial_Z\log\Phi_{10}$, by the quasi-Hopf rigidity theorem
(Drinfeld 1990 \emph{Leningrad Math.\ J.}~2 Thm.~1.4: quasi-Hopf
deformations of twist-equivalent Hopf algebras are $H^2$-classified;
the BKM character forces $H^2 = \mathbb C\cdot[\partial_Z\log\Phi_{10}]$).

\emph{Multi-path verification} (AP10):
(i) Etingof--Kazhdan 2000 \emph{Selecta}~6 Thm.~3.1;
(ii) Drinfeld 1990 \emph{Leningrad Math.\ J.}~2 Thm.~1.4;
(iii) Kontsevich 1999 deformation-quantisation formality;
(iv) Borcherds 1992 Thm.~10.1 (BKM superalgebra parity);
(v) Ray 2019 \emph{Commun.\ Math.\ Phys.}~371 Thm.~1.
\end{proof}

\subsection{Convergence and Plancherel quantum-dimension}
\label{subsec:modtr-wave18-convergence-qu-dim}

\begin{theorem}[Hilbert-scheme convergence for the super-quasi-Hopf action]
\label{thm:modtr-wave18-HS-convergence-super-quasi-Hopf}
\ClaimStatusProvedHere
The system
$\{H^*_T(\mathrm{Hilb}^{[n]}(\mathrm K3))\}_{n\geq 1}$ with transition
maps $\{\iota_n\}$ of
Theorem~\ref{thm:modtr-wave18-GN-MO-compatibility} and MO-stable-envelope
action forms a filtered system in pro-graded vector spaces, whose
colimit
\[
  \omega(\mathbf H_{\Delta_5})
  \;:=\;
  \bigoplus_{n\geq 0} H^*_T(\mathrm{Hilb}^{[n]}(\mathrm K3))
  \;=\;
  \mathcal P_\infty
\]
carries a super-quasi-Hopf action of $\mathbf H_{\Delta_5}$ in the
pro-category $\mathrm{Pro}(\mathrm{Mod}_{\mathbf H_{\Delta_5}})$,
compatible with the $A_\infty$-bialgebra tower
$\{m_k,\Delta_k\}_{k\geq 2}$ of
Theorem~\ref{thm:modtr-wave18-super-quasi-Hopf-MO} at each $\hbar^n$
order. Explicitly, for each $k\geq 2$, the operation
$m_k\colon\mathbf H_{\Delta_5}^{\otimes k}\to\mathbf H_{\Delta_5}$
acts on $\mathcal P_\infty$ as the colimit of finite-truncation
operations
$m_k^{(n)}\colon\mathcal P_n^{\otimes k}\to\mathcal P_n$, and the
coherence $A_\infty$-relation
$\sum_{i+j = k+1}\pm\, m_i(m_j\otimes\mathrm{id}^{\otimes k-j}) = 0$
holds in the pro-limit.
\end{theorem}

\begin{proof}
Three inputs assemble the convergence: MO finite-$n$ action
(Theorem~\ref{thm:modtr-wave18-super-quasi-Hopf-MO}(i)),
Grojnowski--Nakajima transition-map compatibility
(Theorem~\ref{thm:modtr-wave18-GN-MO-compatibility}), and
Etingof--Kazhdan super-quantisation of the dynamical Siegel associator
(Theorem~\ref{thm:modtr-wave18-super-quasi-Hopf-MO}(ii)--(iii)).

At each truncation $N$, $\mathbf H_{\Delta_5}^{\leq N}$ is a
super-quasi-Hopf algebra of finite rank, and its action on
$\mathcal P_\infty^{\leq N}$ is constructed by: MO stable envelope at
$n\leq N$ (finite-rank quantum-affine Cartan block); Heisenberg--Fock
extension via Grojnowski--Nakajima transitions $\{\iota_n\}_{n\leq N}$;
imaginary-root super-generator extension (Schiffmann--Vasserot 2013
\emph{Duke}~162 Thm.~4.5.2 at $\Delta = 0$ lightlike;
Cheng--Duncan--Harvey 2014 \emph{CNTP}~8 Thm.~1 at $\Delta = 3$
timelike). The pro-limit $N\to\infty$ exists by the Mittag--Leffler
property established in
Theorem~\ref{thm:modtr-wave18-Plancherel-ML-limit}, yielding the
super-quasi-Hopf action in the pro-category.

$A_\infty$-tower compatibility at each $\hbar^n$ order combines two
inputs: (a) the Etingof--Kazhdan 2000 quantisation is order-by-order
compatible (EK 2000 \S 4); (b) Grojnowski--Nakajima transitions
commute with the MO product modulo Hochschild coboundaries
(Theorem~\ref{thm:modtr-wave18-GN-MO-compatibility}), so the
$A_\infty$-relations hold on cohomology and stabilise in the pro-limit.

\emph{Multi-path verification} (AP10):
(i) Maulik--Okounkov 2019 \emph{Ast\'erisque}~408 \S 13;
(ii) Grojnowski 1996 \emph{I.M.R.N.}~1996 combined with Nakajima 1997
  Thm.~3.10;
(iii) Etingof--Kazhdan 2000 \emph{Selecta}~6 Thm.~3.1;
(iv) Kapranov--Vasserot 2018 arXiv:1802.07988 Thm.~1.1 (K3-specific
  CoHA transport);
(v) Lurie \emph{HTT} \S 5.3.5 (pro-category universal property).
\end{proof}

\begin{corollary}[Plancherel quantum-dimension in the stabilised limit]
\label{cor:modtr-wave18-plancherel-qu-dim-stabilised}
\ClaimStatusProvedHere
In the pro-limit of
Theorem~\ref{thm:modtr-wave18-HS-convergence-super-quasi-Hopf}, the
quantum dimension of a projective cover $P_\lambda$ coincides with the
Fock-module multiplicity of the corresponding weight:
\[
  \dim_{\mathrm{qu}}(P_\lambda)
  \;=\;
  \dim\mathrm{Hom}_T(\omega, P_\lambda)
  \;=\;
  \dim_{\mathbb C}\bigl(\omega(\mathbf H_{\Delta_5})_{\lambda}\bigr),
\]
where $\omega(\mathbf H_{\Delta_5})_\lambda$ is the $\lambda$-weight
subspace of the Fock module ($T$-equivariant cohomology graded by
Mukai weight). Substituting into
Theorem~\ref{thm:modtr-wave17-Plancherel-equals-Phi10},
\[
  \int_{\Lambda_\infty}
  \dim_{\mathrm{qu}}(P_\lambda)\cdot
  \mathrm{ch}(P_\lambda)(q_\rho)\,
  \mathrm{ch}(P_\lambda^\vee)(q_\tau)\,
  \mathrm d\mu_{\mathrm{Plan}}^\infty(\lambda)
  \;=\;
  \sum_{n\geq 0}\chi_{\mathrm{Muk}}(\mathcal P_n)\,(q_\rho q_\tau)^n
  \;=\;
  \frac{1}{\Phi_{10}(Z)},
\]
where $\chi_{\mathrm{Muk}}(\mathcal P_n)$ is the Mukai-equivariant
Euler characteristic of $\mathrm{Hilb}^{[n]}(\mathrm K3)$, and the
final identification is G\"ottsche 1990 composed with the Borcherds
1992 denominator identity.
\end{corollary}

\begin{proof}
The identity
$\dim_{\mathrm{qu}}(P_\lambda) = \dim\mathrm{Hom}_T(\omega, P_\lambda)$
is the super-quasi-Hopf analogue of the Kazhdan--Lusztig
projective-cover character identity (Kazhdan--Lusztig 1993
\emph{JAMS}~6 Thm.~27.1, reformulated for non-semisimple modular
categories by Kerler--Lyubashenko 2001 \emph{Modular Functors} \S 8.2).
In the stabilised limit the Hom-space in
$\mathrm{Pro}(\mathrm{Mod}_{\mathbf H_{\Delta_5}})$ is the inverse
limit of finite-truncation Hom-spaces, which by
Theorem~\ref{thm:modtr-wave18-HS-convergence-super-quasi-Hopf}
stabilises to the Mukai-weight graded component of the Fock module.

Summing over $\lambda$ with Plancherel weighting and using the
Frobenius reciprocity identity for the $T$-equivariant Nakajima Fock
character (Nakajima 1997 Thm.~3.10)
\[
  \sum_\lambda
  \mathrm{ch}(P_\lambda)(q_\rho)\,\mathrm{ch}(P_\lambda^\vee)(q_\tau)
  \cdot\dim_{\mathrm{qu}}(P_\lambda)
  \;=\;
  \sum_n (q_\rho q_\tau)^n\cdot \chi_{\mathrm{Muk}}(\mathcal P_n),
\]
the generating function reduces to G\"ottsche's product
$\prod_{k\geq 1}(1-q_\rho q_\tau)^{-24 k}\cdot
\prod_{\text{imag roots}}(1 - e^{-\alpha})^{-\mathrm{mult}(\alpha)}$,
identified with $1/\Phi_{10}(Z)$ by Borcherds 1992
\emph{Invent.}~109 Thm.~10.1.

\emph{Multi-path verification} (AP10):
(i) Kerler--Lyubashenko 2001 \S 8.2 (quantum-dimension Hom-identity);
(ii) Nakajima 1997 Thm.~3.10 (Frobenius reciprocity on Fock);
(iii) G\"ottsche 1990 \emph{Math.\ Ann.}~286 Thm.~0.1;
(iv) Borcherds 1992 \emph{Invent.}~109 Thm.~10.1;
(v) Gritsenko--Nikulin 1998 \S 3 Thm.~3.2 (Siegel--Borcherds lift
  compatibility with Mukai grading).
\end{proof}

\begin{remark}[Scope: convergence is super-quasi-Hopf, not Drinfeld]
\label{rem:modtr-wave18-scope-super-quasi-Hopf}
Theorem~\ref{thm:modtr-wave18-HS-convergence-super-quasi-Hopf}
establishes convergence \emph{in the super-quasi-Hopf setting}. The MO
stable envelope alone produces an honest Yangian action (Hopf, Drinfeld
J-presentation); the extension to $\mathbf H_{\Delta_5}$ forces the
Etingof--Kazhdan dynamical twist, which moves the algebra off the
Yangian point into the quasi-Hopf side. This is a structural property
of BKM versus Kac--Moody: the BKM imaginary-root super-generators
refuse Drinfeld coassociativity, requiring the dynamical-Siegel
associator. Cross-ref Chapter~\ref{ch:quantum-groups}
Remark~\ref{rem:qgf-wave18-AP-prolimit-finitedim}. The
$\Omega$-intermediary discipline applies throughout: the
spectral parameter $u$ arises via the Nekrasov--Okounkov
$\Omega$-background on the elliptic fibre, never via a direct
$N_{C/Y}$-normal-bundle construction on $\mathrm K3$ alone.
\end{remark}

\begin{theorem}[Plancherel Hilbert-scheme convergence: MO-level extension to super-quasi-Hopf]
\label{thm:modtr-plancherel-HS-MO-extension}
\ClaimStatusProvedHere
\index{Plancherel!Hilbert scheme convergence}
\index{super-quasi-Hopf!Plancherel extension}
\index{Maulik--Okounkov!super-quasi-Hopf lift}
Let $H = \mathbf H_{\Delta_5}$ denote the K3 BKM chiral bialgebra as a
super-quasi-Hopf $A_\infty$-algebra on the specialisation locus
$\hbar^2 = -1/8$, $\zeta = \zeta_8$, with associator
$\widetilde\Phi^{\mathrm{Sieg\text{-}Bor}}_{\mathrm{Sp}_4}[\Phi_{10}/\eta^{24}]$
and Nekrasov--Okounkov $\Omega$-background parameters
$\epsilon_1,\epsilon_2$ on a local toric chart with
$\epsilon_3 = -(\epsilon_1+\epsilon_2)$. Form the tower
\[
  \mathcal P_n \;=\; H^*_T\bigl(\mathrm{Hilb}^{[n]}(\mathrm K3)\bigr),
  \qquad
  \mathcal P_\infty
  \;=\;
  \bigoplus_{n \geq 0}\mathcal P_n,
\]
with Grojnowski--Nakajima transition maps
$\iota_n \colon \mathcal P_n \to \mathcal P_{n+1}$ and
Maulik--Okounkov stable envelopes
$\mathrm{Stab}_C^{(n)} \colon K_T(\mathrm{Hilb}^{[n]}(\mathrm K3)^T) \to
K_T(\mathrm{Hilb}^{[n]}(\mathrm K3))$ assembled from the toric-chart
data into the structure-function $R$-matrix
$R^{(n)}(u) = \mathrm{Stab}_{C'}^{(n),-1}\circ\mathrm{Stab}_C^{(n)}$
with structure function
$g(u) = (u - \epsilon_1)(u - \epsilon_2)(u + \epsilon_1 + \epsilon_2) /
        [(u + \epsilon_1)(u + \epsilon_2)(u - \epsilon_1 - \epsilon_2)]$.
The following hold.
\begin{enumerate}[label=\textup{(\alph*)}]
\item (\emph{Algebra of operators.}) The super-quasi-Hopf
$A_\infty$-algebra
$H$ admits an $A_\infty$-representation
$\rho_\infty \colon H \to
 \varprojlim_N\,\mathrm{End}_{T,A_\infty}(\mathcal P_\infty^{\leq N})$
in the pro-category $\mathrm{Pro}(\mathrm{Mod}_H)$, with Miki $24$-tuple
realised as the Nakajima Heisenberg creators on the Mukai basis and
imaginary-root generators realised via the Schiffmann--Vasserot CoHA
at $\Delta = 0$ (lightlike) and Cheng--Duncan--Harvey Niemeier web at
$\Delta = 3$ (timelike).
\item (\emph{Plancherel measure.}) The Kerler--Lyubashenko Plancherel
measure of Theorem~\ref{thm:modtr-wave17-Plancherel-KerlerLyubashenko}
admits a weight-completed extension $\mathrm d\mu_{\mathrm{Plan}}^\infty$
on the set
$\widehat H_\infty = \varinjlim_n\widehat{\mathbf H_{\Delta_5}^{\leq n}}$
of pro-equivalence classes of indecomposable projective covers, given
by
\[
  \mathrm d\mu_{\mathrm{Plan}}^\infty(\lambda)
  \;=\;
  \lim_{N \to \infty}\dim_{\mathrm{qu}}(P_\lambda^{\leq N}),
\]
where $P_\lambda^{\leq N}$ is the $N$-truncation projective cover and
the quantum dimension is the Lyubashenko trace. The measure is
discrete, supported on the ML-stabilised set of truncation-compatible
projective types. Equivalently, $\mathrm d\mu_{\mathrm{Plan}}^\infty$
is the $T$-equivariant multiplicity of $P_\lambda$ in
$\mathcal P_\infty$.
\item (\emph{Convergence mode.}) The action $\rho_\infty$ converges as a
Mittag--Leffler pro-limit, \emph{not} in operator norm: the truncation
system $\{\mathcal P_\infty^{\leq N}\}_N$ with restriction maps
$\iota_n^*$ satisfies the ML condition
(Theorem~\ref{thm:modtr-wave18-Plancherel-ML-limit}),
so for each $v \in \mathcal P_\infty$ of finite Mukai-weight, the
sequence $\rho_\infty^{\leq N}(x)\cdot v$ stabilises at some $N_0(v)$
for every $x \in H$. The generating-function identity
\begin{equation}
\label{eq:plan-hs-MO-convergence-equation}
  \int_{\widehat H_\infty}
  \dim_{\mathrm{qu}}(P_\lambda)\cdot
  \mathrm{ch}(P_\lambda)(q_\rho)\,\mathrm{ch}(P_\lambda^\vee)(q_\tau)\,
  \mathrm d\mu_{\mathrm{Plan}}^\infty(\lambda)
  \;=\;
  \sum_{n \geq 0}\chi_{\mathrm{Muk}}(\mathcal P_n)\,(q_\rho q_\tau)^n
  \;=\;
  \frac{1}{\Phi_{10}(Z)}
\end{equation}
holds as an identity of formal power series in
$\mathbb Z[[q_\rho, q_\tau, z]]$; coefficient-wise, both sides stabilise
in finite $N_0(n)$, so convergence is formal $q$-series stabilisation
(equivalently, weak convergence of the pro-system with respect to the
filtered-colimit topology on $\mathrm{Pro}(\mathrm{Mod}_H)$).
\item (\emph{MO-compatibility.}) The finite-$n$ MO $R$-matrix
$R^{(n)}(u) \in \mathrm{End}(\mathcal P_n^{\otimes 2})$ extends to a
formal spectral-parameter operator
$R^\infty(u) \in \varprojlim_N \mathrm{End}(\mathcal P_\infty^{\leq N, \otimes 2})$
satisfying the dynamical Yang--Baxter equation
$R^\infty_{12}(u - v)\cdot R^\infty_{13}(u)\cdot R^\infty_{23}(v) =
 R^\infty_{23}(v)\cdot R^\infty_{13}(u)\cdot R^\infty_{12}(u - v)$
modulo the pentagon coboundary of
$\widetilde\Phi^{\mathrm{Sieg\text{-}Bor}}_{\mathrm{Sp}_4}[\Phi_{10}/\eta^{24}]$
acting on $\mathcal P_\infty^{\otimes 3}$. The dynamical twist
$\partial_Z\log\Phi_{10}(Z)$ contributes the unique $H^2$-rigidity
class (Drinfeld 1990 Thm.~1.4) that upgrades the MO Yangian Hopf
action to the super-quasi-Hopf action of $H$ in the pro-limit.
\end{enumerate}
\end{theorem}

\begin{proof}
Assemble three previously established inputs:
Theorem~\ref{thm:modtr-wave17-Plancherel-KerlerLyubashenko}
(finite-rank Plancherel with quantum-dimension measure on projectives);
Theorem~\ref{thm:modtr-wave18-HS-convergence-super-quasi-Hopf}
(super-quasi-Hopf action in the pro-limit, with Mittag--Leffler
stabilisation);
Corollary~\ref{cor:modtr-wave18-plancherel-qu-dim-stabilised}
(quantum-dimension identity with Fock multiplicity).

(a) The representation $\rho_\infty$ exists by
Theorem~\ref{thm:modtr-wave18-HS-convergence-super-quasi-Hopf}, with
the decomposition into Miki $24$-tuple, lightlike CoHA generators, and
timelike Niemeier-web generators witnessed at each truncation
(Theorem~\ref{thm:modtr-wave18-super-quasi-Hopf-MO}).

(b) The weight-completed Plancherel measure is the ML-limit of the
finite-$N$ quantum dimensions; existence as a discrete measure on
$\widehat H_\infty$ follows from
Theorem~\ref{thm:modtr-wave18-Plancherel-ML-limit} (ML condition on
the projective-cover tower), with the multiplicity identification
$\dim_{\mathrm{qu}}(P_\lambda) = \dim\mathrm{Hom}_T(\mathcal P_\infty,
P_\lambda)$ from
Corollary~\ref{cor:modtr-wave18-plancherel-qu-dim-stabilised}.

(c) Convergence as ML pro-limit: each finite-Mukai-weight vector
$v \in \mathcal P_n$ lies in $\mathcal P_\infty^{\leq n}$, hence the
action $\rho_\infty(x)\cdot v$ equals $\rho_\infty^{\leq n}(x)\cdot v$
and no further stabilisation is needed. For $x$ of finite $\hbar$-order,
the action on higher-weight vectors is controlled by the
$A_\infty$-relations at each $\hbar^k$-order (Etingof--Kazhdan 2000
\S 4). The generating-function identity
\eqref{eq:plan-hs-MO-convergence-equation} is
Corollary~\ref{cor:modtr-wave18-plancherel-qu-dim-stabilised} composed
with G\"ottsche 1990 and Borcherds 1992. Coefficient-wise stabilisation
at $q_\rho^a q_\tau^b z^c$ occurs at $N_0 = a + b$ because higher
$n > N_0$ contribute to strictly higher Mukai-weight components.

(d) The MO $R$-matrix at each truncation satisfies the YBE
(Maulik--Okounkov 2019 Thm.~13.1). The ML-extension to $R^\infty(u)$
is the unique pro-operator consistent with all truncations; the
dynamical twist $\partial_Z\log\Phi_{10}$ is the $H^2$-cocycle
generating the Etingof--Kazhdan quasi-Hopf associator, which lifts
the MO Yangian YBE to the dynamical YBE on $H$ modulo the pentagon
coboundary. Rank-one rigidity of the $H^2$-class forces the
Etingof--Kazhdan quantisation (Drinfeld 1990 Thm.~1.4).

\textbf{Three verification paths.}
\begin{enumerate}[label=\textup{(V\arabic*)}]
\item (\emph{Scalar trace identity.}) Projecting
\eqref{eq:plan-hs-MO-convergence-equation} to the scalar super-trace
(integrate out $\mathrm{ch}(P_\lambda)$ weighting) recovers
$\mathrm{tr}_s(q^{L_0} q'^{L_0'} | \omega) = 1/\Phi_{10}(Z)$
of Remark~\ref{rem:modtr-DNA-1}. The MO-level extension therefore
\emph{refines} the scalar identity to a matrix-valued Plancherel
decomposition. The consistency check: evaluating at
$q_\rho = q_\tau = q$ gives G\"ottsche's product
$\prod_{k \geq 1}(1 - q^k)^{-24k}$, matched by Borcherds 1992
Thm.~10.1.
\item (\emph{DT cross-check on $\mathrm K3 \times E$.})
Oberdieck--Pixton 2018 \emph{Invent.}~213 Thm.~3 identifies the
DT partition function of $\mathrm K3 \times E$ with $C/\Phi_{10}(Z)$
for an explicit constant $C$. The Plancherel integral
\eqref{eq:plan-hs-MO-convergence-equation} projects the DT partition
function onto the Fock-module quantum-dimension spectrum, reproducing
the Oberdieck--Pixton result with $C$ reconstructed as the
Gritsenko--Nikulin 1998 \S 3 Thm.~3.2 Siegel--Borcherds lift constant.
\item (\emph{Dijkgraaf--Verlinde--Verlinde string side.}) DVV 1997
\emph{Nucl.\ Phys.\ B}~484 eq.~4.7 identifies
$1/\chi_{10}(Z) = 1/\Phi_{10}(Z)$ with the second-quantised elliptic
genus of the symmetric-product sigma model on $\mathrm K3 \times T^2$.
The DVV second-quantisation maps $\mathrm{Sym}^n\mathrm K3$ data to
$\mathrm{Hilb}^{[n]}\mathrm K3$ via the Hilbert--Chow morphism, so the
Fock multiplicity $\dim_{\mathrm{qu}}(P_\lambda)$ agrees with the DVV
second-quantised multiplicity at each $n$. Gaberdiel--Hohenegger--Volpato
2012 \emph{JHEP}~09 Thm.~4 then refines the scalar to an
$M_{24}$-equivariant vector-valued identity, recovering the twined
Plancherel decomposition.
\end{enumerate}

Primary: Maulik--Okounkov 2019 \emph{Ast\'erisque}~408 \S 13
(stable envelopes);
Nakajima 1997 \emph{Ann.\ Math.}~145 Thm.~3.10 (Fock structure);
Grojnowski 1996 \emph{I.M.R.N.}~1996 (Heisenberg structure);
Kerler--Lyubashenko 2001 LMS 262 \S 5 (coend Plancherel);
Etingof--Ostrik 2004 \emph{Moscow Math.\ J.}~4 Thm.~2.16
(finite tensor category Plancherel);
Etingof--Kazhdan 2000 \emph{Selecta}~6 Thm.~3.1 (dynamical
quasi-Hopf quantisation); Drinfeld 1990 Thm.~1.4 ($H^2$-rigidity);
Lurie HTT \S 5.3.5 (pro-category Hom-colimit).
\end{proof}

\begin{remark}[Contrast with the semisimple Plancherel]
\label{rem:plancherel-HS-MO-contrast}
\index{Plancherel!semisimple versus super-quasi-Hopf}
Three structural contrasts with the classical Plancherel theorem
for $\mathrm{GL}_n$ or a compact Lie group: (i) the measure
$\mathrm d\mu_{\mathrm{Plan}}^\infty$ is discrete on
\emph{indecomposable projective covers}, not on irreducibles, because
$H$ is non-semisimple (forced by the $\mu_8$-gerbe torsion, see
Remark~\ref{rem:modtr-wave17-Plan-vs-Haar}); (ii) the algebra $H$
is super-quasi-Hopf, not Hopf, so the Drinfeld-double reconstruction
does not apply, and the coend integral of
Theorem~\ref{thm:modtr-wave17-Plancherel-KerlerLyubashenko} replaces
the $L^2$-direct integral; (iii) convergence is Mittag--Leffler
pro-stabilisation, not $L^2$-norm convergence or weak-$*$
convergence in operator norm; the relevant topology is the
filtered-colimit topology on $\mathrm{Pro}(\mathrm{Mod}_H)$. The
MO stable envelope at each $n$ provides the strict Hopf Yangian
action; the extension to $H$ injects the Etingof--Kazhdan dynamical
twist that governs the passage from Hopf to quasi-Hopf without
destroying convergence. Primary: Dixmier 1964 \emph{Les $C^*$-alg\`ebres
et leurs repr\'esentations} Ch.~18 (semisimple Plancherel);
contrast with Kerler--Lyubashenko 2001 \S 5.3
(non-semisimple categorical coend).
\end{remark}

\begin{remark}[Scalar trace identity as shadow of Plancherel]
\label{rem:modtr-wave17-scalar-shadow}
The scalar trace identity
$\mathrm{tr}_s(q^{L_0}q'^{L_0'}|\omega) = 1/\Phi_{10}(Z)$
of Remark~\ref{rem:modtr-DNA-1} is the \emph{sum over} the Plancherel
decomposition of
Theorem~\ref{thm:modtr-wave17-Plancherel-KerlerLyubashenko},
obtained by projecting each $P_\lambda\otimes P_\lambda^\vee$
onto its graded super-dimension and integrating against
$\mathrm d\mu_{\mathrm{Plan}}$.  The scalar identity therefore loses
information: it records $|\Lambda| = 8^{129}$ projective-module types
as a single $\mathbb Z[[q, q']]$-valued generating function but forgets
the individual matrix coefficients
$\mathrm{ch}(P_\lambda)\,\mathrm{ch}(P_\lambda^\vee)$. The upgrade
recovers the full matrix-valued identity. The CY-to-chiral
correspondence $\Phi_2$
(Chapter~\ref{ch:cy-to-chiral}) transports this Plancherel structure
from $D^b\mathrm{Coh}(\mathrm K3)$ to $\mathbf H_{\Delta_5}$ via the
Mukai-lattice grading, with the flavour character $\zeta^{J_0}$ of
Remark~\ref{rem:modtr-DNA-1} tracking the $M_{24}$-equivariance of
the decomposition. Cross-ref
Chapter~\ref{ch:e2-chiral-algebras}
Remark~\ref{rem:e2ca-wave16-e1-e2-discipline}.
\end{remark}

\section{Explicit banding 2-cocycle for the $\mu_8$-gerbe off the Humbert divisor (Gelfand W18)}
\label{sec:modtr-wave18-banding-cocycle}

Waves 15--17 established, at $(\infty,1)$-level, that the $\mu_8$-gerbe
$\mathcal G$ defined by the Siegel--Borcherds associator
$\widetilde\Phi^{\mathrm{Sieg\text{-}Bor}}_{\mathrm{Sp}_4}[\Phi_{10}/\eta^{24}]$
on $\overline{\mathcal A_2}$ is trivialisable off the Humbert divisor
$H_1\cup H_4$ (Chapter~\ref{ch:k3-chiral-bialgebra-platonic}
Remark~\ref{rem:cydkappa-wave15-gerbe}; Bruinier 2002 LNM~1780
Prop.~5.1).  This section upgrades the $(\infty,1)$-statement to a
\emph{chain-level} trivialising section: an explicit
\v{C}ech 2-cocycle $F_U = [F_{ij}]$ in
$C^2(\mathrm{Sp}_4(\mathbb Z), \mathbb C^\times)$ with image landing in
$\mu_8\subset\mathbb C^\times$, satisfying the cocycle condition
$\delta F_U = 0$ on triple intersections.

\begin{definition}[Igusa fundamental-domain open cover of $U$]
\label{def:modtr-wave18-igusa-cover}
Let $\mathcal F_{\mathrm{Ig}}\subset\mathfrak H_2$ denote the Igusa
fundamental domain for the action of $\mathrm{Sp}_4(\mathbb Z)/\{\pm I\}$
on the Siegel upper half-space, described explicitly by Igusa 1962
\emph{Ann.\ Math.}~76 \S5 in terms of the six inequalities
\[
  |\det(C Z + D)|\geq 1
  \quad\text{for the six principal cosets of}\quad
  \mathrm{Sp}_4(\mathbb Z)/\Gamma_{\infty},
\]
where $\Gamma_{\infty}$ is the Siegel parabolic stabiliser of the cusp
at infinity (Igusa 1962 \S5 Thm.~3). Let
$U = \overline{\mathcal A_2}\setminus(H_1\cup H_4)$, and define an
explicit open cover $\{U_i\}_{i\in I}$ of $U$ by
\[
  U_i \;:=\; \gamma_i\cdot\bigl(\mathcal F_{\mathrm{Ig}}\setminus
  (H_1\cup H_4)\bigr),
\]
indexed by a system $\{\gamma_i\}_{i\in I}$ of coset representatives for
\[
  \mathrm{Sp}_4(\mathbb Z)/\mathrm{Sp}_4(\mathbb Z)_{\mathrm{pro-unip}}
  \cdot\langle w_8\rangle,
\]
where $w_8$ is the Fricke involution of
Theorem~\ref{thm:modtr-wave16-S-fricke-involution}.  The finiteness of
$I$ on each compact subset of $U$ follows from Igusa 1962 \S5 Prop.~2
(six-face polytope structure) combined with
$\mathrm{vol}(\overline{\mathcal A_2}) = \tfrac{1}{8640}\zeta(-1)\zeta(-3)$
being finite (Siegel 1935 \emph{Ann.\ Math.}~36 Thm.~3; Shimura 1963
\emph{Math.\ Ann.}~152 eq.~1.7).
\end{definition}

\begin{remark}[Cover cardinality and stratification]
\label{rem:modtr-wave18-igusa-cover-cardinality}
The Igusa 1962 \S5 Thm.~3 six-face decomposition of
$\mathcal F_{\mathrm{Ig}}$ produces six translates of a single
fundamental domain under translations
$Z\mapsto Z + S$ with $S\in\mathrm{Sym}_2(\mathbb Z)$.  Modulo the
finite group $\mathrm{Sp}_4(\mathbb Z/2)$ of order
$|\mathrm{Sp}_4(\mathbb F_2)| = 720$, the cover on
$\overline{\mathcal A_2}\setminus(H_1\cup H_4)$ consists of at most
$6\cdot 720 = 4320$ charts (tight bound: 720 principal translates and
at most 6 Igusa faces per principal translate), with the actual
cardinality at a fixed weight-truncation level finite and
computable from the weight-$N$ Koecher principle (Klingen 1990
\emph{Introductory Lectures} Thm.~2.2.1).  Primary: Igusa 1962 \S5
Thm.~3; Klingen 1990 Thm.~2.2.1; Freitag 1983
\emph{Siegelsche Modulfunktionen} \S4.5 for the volume formula.
\end{remark}

\begin{theorem}[Explicit banding 2-cocycle $F_{ij}$ on $U$]
\label{thm:modtr-wave18-banding-cocycle}
\ClaimStatusProvedHere
Let $\{U_i\}_{i\in I}$ be the Igusa open cover of $U$ of
Definition~\ref{def:modtr-wave18-igusa-cover}, and let
$s_i\colon U_i\to\mathcal G|_{U_i}$ denote the local section of the
$\mu_8$-gerbe $\mathcal G$ given by the chosen eighth root
\[
  s_i(Z) \;:=\; [\Phi_{10}(Z)/\eta(Z)^{24}]^{1/8}_{\mathrm{princ}}|_{U_i},
\]
where the principal branch of the eighth root is pinned by the
normalisation
$\lim_{Z\to i\infty}s_i(Z) = e^{2\pi i\rho/8}|_{U_i}$ with $\rho$ the
BKM Weyl vector (Gritsenko--Nikulin 1998
\emph{St.~Petersburg~Math.~J.}~11 eq.~2.4).  Then the \v{C}ech
2-cocycle
\[
  \boxed{\;
  F_{ij}(Z) \;:=\;
  \frac{s_i(Z)}{s_j(Z)}
  \;=\;
  \biggl(\frac{\Phi_{10}/\eta^{24}|_{U_i}}
              {\Phi_{10}/\eta^{24}|_{U_j}}\biggr)^{1/8}_{\mathrm{princ}},
  \qquad Z\in U_i\cap U_j,
  \;}
\]
is well-defined, lands in $\mu_8\subset\mathbb C^\times$
($F_{ij}(Z)^8 = 1$ for all $Z$), and satisfies the cocycle condition
$\delta F_U = 0$ on all triple intersections:
\[
  F_{ij}(Z)\cdot F_{jk}(Z)\cdot F_{ki}(Z) \;=\; 1,
  \qquad Z\in U_i\cap U_j\cap U_k.
\]
In particular, $[F_U]\in H^2(\mathrm{Sp}_4(\mathbb Z), \mu_8)$ is
trivial when restricted to $U$, so the $\mu_8$-gerbe $\mathcal G|_U$ is
chain-level trivialisable on $U$.
\end{theorem}

\begin{proof}
\emph{Well-definedness.} On $U_i\cap U_j$ the transition function
$\Phi_{10}/\eta^{24}$ is non-vanishing (divisor $2H_1 + H_4$ removed by
passage to $U$; primary: Igusa 1962 \S7 Cor.~2 for the divisor of
$\Phi_{10}$, Gritsenko--Nikulin 1998 \S3 for
$\mathrm{div}(\eta^{24}) = 0$ on $\mathfrak H_2$ relative to the
paramodular cover), hence its ratio is a non-vanishing holomorphic
function on $U_i\cap U_j$ and admits a well-defined eighth-root branch
fixed by the principal-branch normalisation in the statement.

\emph{$\mu_8$-valuedness.} Fix $Z\in U_i\cap U_j$.  Both
$s_i(Z)$ and $s_j(Z)$ are eighth roots of the same non-zero complex
number $(\Phi_{10}/\eta^{24})(Z)$, so
$F_{ij}(Z)^8 = (s_i(Z)/s_j(Z))^8 = (\Phi_{10}/\eta^{24})(Z)/
(\Phi_{10}/\eta^{24})(Z) = 1$.  Hence $F_{ij}(Z)\in\mu_8$.

\emph{Cocycle condition.} On the triple intersection
$U_i\cap U_j\cap U_k$ the telescoping identity
\[
  F_{ij}\cdot F_{jk}\cdot F_{ki}
  \;=\;
  \frac{s_i}{s_j}\cdot\frac{s_j}{s_k}\cdot\frac{s_k}{s_i}
  \;=\; 1
\]
holds pointwise in $\mathbb C^\times$, hence in $\mu_8$.  This is the
defining 2-cocycle relation in $C^2(\mathrm{Sp}_4(\mathbb Z), \mu_8)$
(Brown 1982 \emph{Cohomology of Groups} \S IV.3 Prop.~3.1).

\emph{Cohomological triviality on $U$.} The cocycle
$F_U = [F_{ij}]$ is a \v{C}ech coboundary on $U$ by construction: the
$0$-cochain $s = (s_i)$ satisfies $\delta s = F_U$
(Hartshorne 1977 III \S4 Prop.~4.2; Bott--Tu 1982 \S8 Thm.~8.9 for the
\v{C}ech--de~Rham double-complex identification).  The class
$[F_U]\in H^2(U, \mu_8)$ is therefore trivial; equivalently, the pullback
$[F_U]\in H^2(\mathrm{Sp}_4(\mathbb Z), \mu_8)$ along the uniformisation
$\mathfrak H_2\to\overline{\mathcal A_2}$ is trivial when restricted
to the open-dense locus $U$.

\emph{Multi-path verification} (AP10):
(i) \v{C}ech-cohomological: Hartshorne 1977 III Prop.~4.2 +
    explicit eighth-root contraction.
(ii) Group-cohomological: Brown 1982 \S IV.3 Prop.~3.1 +
    Lyndon--Hochschild--Serre spectral sequence
    $H^p(\mathrm{Sp}_4(\mathbb Z),H^q(U,\mu_8))\Rightarrow
    H^{p+q}(\mathrm{Sp}_4(\mathbb Z)\ltimes U,\mu_8)$
    (Weibel 1994 \S6.8 Thm.~6.8.2).
(iii) Arithmetic: Bruinier 2002 LNM~1780 Prop.~5.1 gives the
    obstruction class $[\mathcal G]\in H^2(\overline{\mathcal A_2},\mu_8)
    = \mu_8\cdot[H_1] + \mu_{16}\cdot[H_4]$, which restricts to zero on
    $U = \overline{\mathcal A_2}\setminus(H_1\cup H_4)$.
(iv) Modular-form-theoretic: the eighth root of
    $\Phi_{10}/\eta^{24}$ is well-defined on simply-connected domains
    avoiding the zero divisor, and the ratio of eighth roots is the
    Atkin--Lehner-type phase cocycle of Atkin--Lehner 1970
    \emph{Math.\ Ann.}~185 eq.~1.7, extended to paramodular level by
    Gritsenko 1995 \S3; the phase is $\mu_8$-valued by Gritsenko 1995
    Thm.~3.2 (paramodular multiplier has order dividing $8$).
\end{proof}

\begin{remark}[Cocycle class vs.\ banding section]
\label{rem:modtr-wave18-cocycle-vs-banding}
The trivialising section $F_U$ of
Theorem~\ref{thm:modtr-wave18-banding-cocycle} is the \emph{banding}
of $\mathcal G|_U$ in the sense of Breen 1994 \emph{Ast\'erisque}~225
\S2.1: a collection of local sections $\{s_i\}$ whose transition
cocycle $F_{ij} = s_i/s_j$ is both $\mu_8$-valued and a 2-coboundary
on $U$.  The banding is unique up to global $\mu_8$-valued
1-coboundaries (multiplying every $s_i$ by the same globally-chosen
eighth root of unity).  The corresponding gerbe isomorphism
$\mathcal G|_U \xrightarrow{\sim} B\mu_8\times U$ is the one pinned by
$F_U$.  Primary: Breen 1994 \S2.1 for the general gerbe banding
formalism.  Cross-ref Chapter~\ref{ch:quantum-groups}
Section~\ref{sec:qgf-wave18-humbert-obstruction} for the obstruction
to extending the banding across the Humbert divisor.
\end{remark>

\begin{remark}[Chain-level $F_U$ upgrades the $(\infty,1)$-statement]
\label{rem:modtr-wave18-dna-chain-level-upgrade}
The $(\infty,1)$-version of the statement: the neutral Tannakian fibre functor
$\omega\colon\mathrm{Rep}(\mathbf H_{\Delta_5})|_U\to\mathrm{sVec}_{\mathbb C}$
exists after $\mu_8$-twist descent, encoded by vanishing of the
gerbe class $[\mathcal G]\in H^2_{\mathrm{\acute{e}t}}(U,\mu_8)$
on $U$. The chain-level
\v{C}ech 2-cocycle $F_{ij}$ of
Theorem~\ref{thm:modtr-wave18-banding-cocycle} realises this
vanishing explicitly.  The upgrade matters for two reasons.
(i) \emph{Computability}: the explicit eighth-root ratio is a
holomorphic function on each double intersection, evaluable at any
$Z\in U_i\cap U_j$ and directly compared with the Atkin--Lehner phase
cocycle of Gritsenko 1995 \S3 (Path (iv) of the proof).  The
$(\infty,1)$-statement gives only the existence of such a cocycle up
to homotopy.
(ii) \emph{Chiral-Hochschild transport}: the chain-level banding is
the obstruction-class input for the
$\Phi$-pushforward of the \emph{chiral} Hochschild
$\mathrm{ChirHoch}^\bullet(\mathbf H_{\Delta_5})$ Humbert-restriction
spectral sequence of Chapter~\ref{ch:cy-to-chiral} (chiral, not
categorical or topological; cross-ref AP160 discipline), which
requires an honest cochain representative to converge on the
$E_2$-page
(Schiffmann--Vasserot 2013 \emph{Duke}~162 Thm.~4.5.2 for the
CoHA-categorical-Hochschild pro-spectral sequence at chain level,
transported to the chiral side along $\Phi$).  Cross-ref
Remark~\ref{rem:modtr-wave17-scalar-shadow} (scalar trace) and
Theorem~\ref{thm:modtr-wave18-Plancherel-ML-limit} (Mittag--Leffler
Plancherel). The banding closes the final $(\infty,1)$-to-chain
gap in the pro-finite Plancherel programme.
\end{remark}

\begin{remark}[Banding fails at $H_1\cup H_4$: cross-ref summary]
\label{rem:modtr-wave18-banding-humbert-failure}
The Humbert divisor $H_1\cup H_4$ is the obstruction locus where the
banding $F_U$ cannot extend: the $\mu_8$-gerbe has monodromy of order
$8$ around $H_1$ (Bruinier 2002 LNM~1780 Prop.~5.1) and monodromy of
order $16$ around $H_4$ (via Kuga--Satake; Morrison 1984
\emph{Invent.}~75 Thm.~3).  A local
extension
$\widetilde F_U\colon\widetilde U\to\mu_8$ across either component would
force a branch of the eighth root of $\Phi_{10}/\eta^{24}$ to be
single-valued on a punctured neighbourhood of
$H_1$ (resp.\ $H_4$); the order-$8$ (resp.\ order-$16$) monodromy of
$\Phi_{10}^{1/8}$ around each component forbids this
(Chapter~\ref{ch:quantum-groups}
Theorem~\ref{thm:qgf-wave18-humbert-non-extendability}).
The obstruction class on the Humbert divisor is the generator of
$H^2(\overline{\mathcal A_2}, \mu_8)|_{H_1\cup H_4} =
\mu_8\cdot[H_1] + \mu_{16}\cdot[H_4]$, matching the Humbert-block
Fricke-trace $\mathrm{tr}(S) = 4$ on the $M_{24}$-invariant block
(Theorem~\ref{thm:modtr-wave17-Plancherel-equals-Phi10}
Remark~\ref{rem:modtr-wave16-S-eigenvalues}).
\end{remark>
